\chapter{Variational Iterative Methods for scattering problems } 
\label{variational}
%%%%%%%%%%%%%%%%%%%%%%%%%%%%%%%%%%%%%%%%%%%%%%%%%%%%%%%%%%%%%%%%%%%%%%%%%%%%%%%%%%%%%%%%%%%%


\section{Introduction}

We recall from chapter \ref{stratG}, eqn.(\ref{LSstrat}), that the 
$6\times 6$ 
displacement/traction, Green's tensor field 
${\CB G}$  of a heterogeneous medium,   may be constructed using the 
formula
\begin{equation} \label{LSstratVar}
{\CB G}=\rf{\CB G}+
\rf{\CB P}\!{}^{z,R}_{s} 
{\mit \Delta}{\CB C}{}_{eff}(k_y)
\left({\bmi I} -  
\rf{\CB P}_{r,s}
{\mit \Delta}{\CB C}{}_{eff}(k_y)\right)^{-1} 
\rf{\CB G}\!{}^{z}_{r}
\end{equation}
where the $6\times 6$ array $\rf{\CB G}$ is the corresponding Green's 
tensor for the stratified 
reference component of the model. The $8 L N_x \times 6$ incident field 
$\rf{\CB G}\!{}^{z}_{r}$ was given by 
either (\ref{Gincdown}) or (\ref{Gincup}) depending on whether the the 
source is above 
or below the zone of heterogeneity, respectively. The  $8 L N_x \times 8 L 
N_x $ propagator
$\rf{\CB P}_{r,s}$ was given by (\ref{CBPdef0}).
 This operator maps the scattered field from one scattering interaction, 
to 
the  incident field at the subsequent scattering interaction.
The scattering interactions are governed by the $L N_x \times L N_x$ 
diagonal array
${\mit \Delta}{\CB C}{}_{eff}(k_y)$ of (\ref{dCBC}). Each 
element on the diagonal of ${\mit \Delta}{\CB C}{}_{eff}(k_y)$ is  an 
$8 \times 8$ array that models the scattering interaction at a given point 
in space. 
Finally, the $6 \times 8 L N_x $ propagator
$\rf{\CB P}\!{}^{z,R}_{s}$ carries the scattered field to the 
receiver. This operator was defined in (\ref{rfCBPzRsl}). A schematic 
representation of (\ref{LSstratVar}) is depicted by Fig.\ref{scatGeom}.
%%%%%%%%%%%%%%%%%%%%%%%%%%%%%%%%%%%%%%%%%%%%%%%%%%%%%%%%%%%%%%%%%%%%%%%%%%%%%%%%%%%%%%%%%%%%

%%%%%%%%%%%%%%%%%%%%%%%%%%%%%%%%%%%%%%%%%%%%%%%%%%%%%%%%%%%%%%%%%%%%%%%%%%%%%%%%%%%%%%%%%%%%

%%%%%%%%%%%%%%%%%%%%%%%%%%%%%%%%%%%%%%%%%%%%%%%%%%%%%%%%%%%%%%%%%%%%%%%%%%%%%%%%%%%%%%%%%%%%

\begin{figure}
\HideDisplacementBoxes
%\ShowDisplacementBoxes
\centerline{\BoxedEPSF{scatGeom.EPSF scaled 1000}}
\caption{The canonical scattering geometry for a source and receiver in 
the surface zone $z_f \le z \le z_0$}
\protect\label{scatGeom}
\end{figure}
%%%%%%%%%%%%%%%%%%%%%%%%%%%%%%%%%%%%%%%%%%%%%%%%%%%%%%%%%%%%%%%%%%%%%%%%%%%%%%%%%%%%%%%%%%%%

%%%%%%%%%%%%%%%%%%%%%%%%%%%%%%%%%%%%%%%%%%%%%%%%%%%%%%%%%%%%%%%%%%%%%%%%%%%%%%%%%%%%%%%%%%%%

%%%%%%%%%%%%%%%%%%%%%%%%%%%%%%%%%%%%%%%%%%%%%%%%%%%%%%%%%%%%%%%%%%%%%%%%%%%%%%%%%%%%%%%%%%%%


The task for this chapter, is to derive some approximate constructions 
for ${\CB G}$ from  (\ref{LSstratVar}), while avoiding a direct 
construction of the  enormous matrix inverse 
\begin{equation} \label{moller}
{\bmi \Omega}\equiv
\left({\bmi I} -  
\rf{\CB P}_{r,s}
{\mit \Delta}{\CB C}{}_{eff}(k_y)\right)^{-1}
\end{equation}
in (\ref{LSstratVar}). We note in passing that 
${\bmi \Omega}$ is the time harmonic, elastic wave equation analogue 
of the {\em 
M{\o}ller wave operator\/} of quantum scattering theory 
[e.g. {\em Taylor \/}~(1972)].

We will show how to avoid constructing ${\bmi \Omega}$ using expensive 
matrix inversion,  by employing what amounts to a {\em matrix polynomial 
approximation\/} in
the Lippmann-Schwinger type operator ${\CB H}(k_y)$, 
defined by
\begin{equation} \label{Hdef}
	 {\CB H}(k_y)=\left({\bmi I} -  
\rf{\CB P}_{r,s}(k_y)
{\mit \Delta}{\CB C}{}_{eff}(k_y)\right),
\end{equation}
in terms of the explicitly available matrix operators $\rf{\CB 
P}_{r,s}(k_y)$ and ${\mit \Delta}{\CB C}{}_{eff}(k_y)$ that we have met 
previously. We should point out from the outset, that although the methods 
to be 
developed in this chapter are several orders of magnitude faster than 
direct construction of ${\bmi \Omega}$ in (\ref{moller}) using matrix 
inversion
(at least for large problems), the methods are inherently approximate.
A nice introduction to {\em matrix polynomial constructions of a matrix 
inverse\/} has been given by {\em Vorobyev\/}~(1965). 
 

We employ two classes of methods involving matrix polynomials in ${\CB 
H}(k_y)$:
\begin{itemize}
\item the first method, is a
generalization of the pioneering work of {\em Lapwood \& 
Hudson\/}~(1975) concerning the {\em Schwinger-Levine variational 
formulation of the seismic scattering transition\/}. We generalize their 
work in two respects:
\begin{enumerate}
\item the method is extended to laterally heterogeneous media
\item the basic method is interpreted as the leading order term in a 
recursive sequence of additive corrections. The higher order corrections 
allow us to extend the basic method to {\em strong scattering 
environments\/} 
\end{enumerate}

\item the second method is a generalization of the {\em generalized 
minimal residual method\/} (GMRES) of {\em Saad \& Schultz\/}~(1986).

\end{itemize}

Each of these methods constructs the best matrix polynomial approximation 
in ${\CB H}$ to ${\bmi \Omega}$,  by choosing the undetermined multipliers
in a  polynomial of fixed degree,  so that a {\em variational functional 
\/} is stationary
with respect to small changes in these parameters.
The variational formulation of problems governed by integral equations 
have long been known (Volterra, 1884), and general techniques for the 
construction of variational principles are now well developed (Gerjouy, 
Rau \& Spruch, 1983). 

The symmetry of the field equations (\ref{Hdef}) plays a central role in 
the variational formulation of the problem.
The propagator $\rf{\CB P}_{r,s}(k_y)$ and scattering matrix 
${\mit \Delta}{\CB C}_{eff}(k_y)$  have the following symmetry properties
\begin{equation} \label{sym}
	{\mit \Delta}{\CB C}_{eff}(k_y)\ =\  {\mit \Delta}{\CB 
	C}{}^{T}_{eff}(-k_y) \qquad 
\rf{\CB P}_{r,s}(k_y)\ =\  \rf{\CB P}\!\!{}_{r,s}^{ T}(-k_y)
\end{equation}
Next we introduce a {\em bilinear form\/} $\left\langle{ \cdot \ ,\ \cdot 
}\right\rangle:
{\Bbb{C}}^{8 N_x L}\times {\Bbb{C}}^{8 N_x L} \rightarrow 
{\Bbb{C}}$, defined as 
\begin{equation} \label{bilinearDef}
	\left\langle{ {\bmi  v}_{1} , {\bmi  v}_2 }\right\rangle \equiv {\bmi  
	v}^{T}_{1}   {\mit \Delta}{\CB C}_{eff}(k_y)   {\bmi  v}_2.
\end{equation}  
Unfortunately, $\left\langle{ {\bmi  v}_{1}  , {\bmi  v}_{1} 
}\right\rangle=0 
\not\Rightarrow {\bmi  v}_{1}\equiv{\bf  0}$ since ${\bmi  v}_1$ is 
complex 
valued and 
the weight matrix $ {\mit \Delta}{\CB C}_{eff}(k_y)$ 
used in the definition (\ref{bilinearDef}) is not {\em positive 
definite}.
Using (\ref{sym}) it follows that
\[
{\mit \Delta}{\CB C}_{eff}(k_y) {\left({ {\bmi { I}} -\rf{\CB P}_{r,s}(k_y)
{\mit \Delta}{\CB C}_{eff}(k_y)  }\right)} = {\left({  {\bmi { I}} 
-\rf{\CB P}_{r,s}(-k_y)
{\mit \Delta}{\CB C}_{eff}(-k_y)  }\right)}^{T} {\mit \Delta}{\CB 
C}_{eff}(k_y)
\]
so that
\begin{equation} \label{HsymDef}
	\left\langle{{\CB H}(-k_y) \cdot  {\bmi  v}_1,{\bmi  v}_2}\right\rangle=
	\left\langle{{\bmi  v}_1,{\CB H}(k_y) \cdot {\bmi  v}_2}\right\rangle
\end{equation}
for arbitrary fields ${\bmi  v}_1$ and ${\bmi  v}_2$, i.e. the {\em 
operator 
adjoint\/} ${\CB 
H}^{\dagger}$ of ${\CB H}$ is defined as
\begin{equation} \label{Htdef}
{\CB H}^{\dagger}(k_y)
 \equiv 
{\left({ \ {\bmi { I}} -\rf{\CB P}_{r,s}(-k_y)
	{\mit \Delta}{\CB C}_{eff}(-k_y) \ }\right)}\equiv {\CB H}(-k_y)
\end{equation} 
relative to the bilinear form of  (\ref{bilinearDef}). 


%%%%%%%%%%%%%%%%%%%%%%%%%%%%%%%%%%%%%%%%%%%%%%%%%%%%%%%%%%%%%%%%%%%%%%%%%%%%%%%%%%%%%%%%%%%%%

%%%%%%%%%%%%%%%%%%%%%%%%%%%%%%%%%%%%%%%%%%%%%%%%%%%%%%%%%%%%%%%%%%%%%%%%%%%%%%%%%%%%%%%%%%%%%

%%%%%%%%%%%%%%%%%%%%%%%%%%%%%%%%%%%%%%%%%%%%%%%%%%%%%%%%%%%%%%%%%%%%%%%%%%%%%%%%%%%%%%%%%%%%%


Using the bilinear functional $\left\langle{ \cdot \ ,\ \cdot 
}\right\rangle$ of 
(\ref{bilinearDef}), it follows from (\ref{LSstratVar}) that the 
$(\omega,k_y)$-response of a $3-$D  medium with $2-$D superimposed
 heterogeneity driven by a point force ${\bmi 
 f}(\omega,k_y)$ at ${\bmi  x}_S$ and recorded by a point receiver at  
${\bmi  x}_R$ may
be written as
\begin{eqnarray} \label{response3}
\tilde{\bmi f}\cdot{\CB G}\cdot{\bmi f}&=&
\tilde{\bmi f}\cdot\rf{\CB G}\cdot{\bmi f}+
\tilde{\bmi f}\cdot\rf{\CB P}\!{}^{z,R}_{s} 
{\mit \Delta}{\CB C}{}_{eff}(k_y)
\left({\bmi I} -  
\rf{\CB P}_{r,s}
{\mit \Delta}{\CB C}{}_{eff}(k_y)\right)^{-1} 
\rf{\CB G}\!{}^{z}_{r}\cdot{\bmi f}\nonumber\\
&=&\tilde{\bmi f}\cdot\rf{\CB G}\cdot{\bmi f}+ 
\left\langle\tilde{\bmi b} , {\bmi  \Omega} (k_y)\cdot {\bmi b} 
\right\rangle
\end{eqnarray}
where
\begin{equation} \label{bvardef}
{\bmi b}= \rf{\CB G}\!{}^{z}_{r}\cdot{\bmi f}
\end{equation}
and
\begin{equation} \label{btdef}
\tilde{\bmi b}= \left[\rf{\CB P}\!{}^{z,R}_{s}(k_y)\right]^{T}\cdot 
\tilde{\bmi f}.
\end{equation}
The {\em receiver directivity vector\/} $\tilde{\bmi f}$ 
in (\ref{response3}) picks out the relevant displacement component to be 
recorded at ${\bmi 
x}_R$.
Eqn.(\ref{response3}) decomposes the {\em total\/} response $\tilde{\bmi 
f}\cdot{\CB G}\cdot{\bmi f}$ at a fixed 
receiver position ${\bmi x}_R$ into the sum of a {\em reference\/} 
component
\begin{equation} \label{refcomponent}
\tilde{\bmi f}\cdot\rf{\CB G}\cdot{\bmi f}
\end{equation} 
propagating exclusively through the stratified background medium and a 
{\em scattering transition\/} 
\begin{equation} \label{scattransition}
\left\langle\tilde{\bmi b} , {\bmi  \Omega} 
\cdot {\bmi b} \right\rangle
\end{equation}
induced by the superimposed heterogeneity. 
%%%%%%%%%%%%%%%%%%%%%%%%%%%%%%%%%%%%%%%%%%%%%%%%%%%%%%%%%%%%%%%%%%%%%%%%%
In {\em Tromp \& Dahlen\/}~(1990), an equation for the scattering 
transition analogous to (\ref{response3}) 
was written down to model the {\em normal modes\/} in a heterogeneous 
earth, although the method of derivation, and the intended application 
were 
different from this work.
%%%%%%%%%%%%%%%%%%%%%%%%%%%%%%%%%%%%%%%%%%%%%%%%%%%%%%%%%%%%%%%%%%%%%%%%%
%%%%%%%%%%%%%%%%%%%%%%%%%%%%%%%%%%%%%%%%%%%%%%%%%%%%%%%%%%%%%%%%%%%%%%%%%%%%%%%%%%%%%%%%%

%%%%%%%%%%%%%%%%%%%%%%%%%%%%%%%%%%%%%%%%%%%%%%%%%%%%%%%%%%%%%%%%%%%%%%%%%%%%%%%%%%%%%%%%%

%%%%%%%%%%%%%%%%%%%%%%%%%%%%%%%%%%%%%%%%%%%%%%%%%%%%%%%%%%%%%%%%%%%%%%%%%%%%%%%%%%%%%%%%%

%%%%%%%%%%%%%%%%%%%%%%%%%%%%%%%%%%%%%%%%%%%%%%%%%%%%%%%%%%%%%%%%%%%%%%%%%%%%%%%%%%%%%%%%%


The calculation of the reference component (\ref{refcomponent}) using the 
{\em reflectivity method\/} is now 
routine (Fuchs \& M\"{u}ller, 1971; Kennett, 1983; M\"{u}ller, 1985; 
Mallick \& Frazer, 
1987). In this chapter we will discuss some methods that are effective 
for the calculation of scattered field from the superimposed lateral 
heterogeneity. 
In particular, we consider the following  two strategies to construct the 
{\em scattering 
transition\/} (\ref{scattransition}):
\begin{enumerate}
\item given the incident field (\ref{bvardef}) and the Lippmann-Schwinger 
type operator (\ref{Hdef}), we solve 
\begin{equation} \label{goveqn}
{\CB H}(k_y){\bmi  v} = {\bmi  b},
\end{equation}
for the field vector ${\bmi  v}$ 
using an iterative scheme that is effective for the solution of 
{\em complex valued, indefinite\/} linear systems. Then we form
$\langle \tilde{\bmi  b}, {\bmi  v} \rangle$, which is the required 
scattering transition


\item we construct {\em high order Schwinger-Levine variational 
approximations\/} for a 
required {\em scattering transition\/} $\langle \tilde{\bmi  b}, {\bmi 
\Omega} {\bmi  b} \rangle$
directly. In addition to the {\em primal \/} problem (\ref{goveqn}), 
we introduce the {\em reciprocal\/} or {\em dual\/} problem
\begin{equation} \label{goveqnT}
{\CB H}(-k_y) \cdot \tilde{\bmi v} =\tilde{\bmi b},
\end{equation}
where $\tilde{\bmi b}({\bmi x})$ was given by (\ref{btdef}).

\end{enumerate}

Clearly, once we have solved ${\CB H}{\bmi  v} = {\bmi  b}$, for the field 
vector ${\bmi  v}$
as in class ({\bf 1}), we can construct any required scattering 
transition  
$\langle \tilde{\bmi  b},{\CB H}^{-1} {\bmi  b}\rangle=
\langle \tilde{\bmi  b}, {\bmi  v}\rangle$ in class ({\bf 2}) immediately. 
However, provided the number of required transitions is small (say 
$1-2$), it is usually more efficient in terms of both computation time and 
storage, to solve for the transition directly, without first solving for 
$\bmi  v$. 

All the schemes considered share a common feature:\newline
at iteration $n$, the current approximation ${\bmi  v}_{n}$ to the 
solution 
of ${\CB H} {\bmi  v}= {{\bmi { b}}} $
is chosen to make a {\em variational functional\/} $J({\bmi  v} )$ 
stationary 
(i.e. $\delta J({\bmi  v}_{n} )=0$), where the variational trial 
solution ${\bmi  v}_{n}$ is an element of the {\em trial space\/} defined 
as
\begin{equation} \label{krylov}
{\bmi  v }_{n} \in  {\bmi  v}_{0} 
+{\cal K}_{n}\left({\bmi  r}_{0},{\CB H}\right) ,\ \ \ \ \ n=1,2,\ldots  .
\end{equation}
with ${\bmi  v}_{0}$ the initial approximation to 
$  {\bmi  v}= {\CB H}^{-1}{{\bmi { b}}} $ , 
${\bmi  r}_{0}= {{\bmi { b}}}- {\CB H}{\bmi  v}_{0}$ 
the corresponding residual vector, and 
\begin{equation} 	\label{Krylovdef}
	{\cal K}_{n}({\bmi  r}_{0},{\CB H}):= \mbox {span}\left\{{{\bmi  
r}_{0},{\CB H}{\bmi  r}_{0},
	\cdots ,{\CB H}^{n-1}{\bmi  r}_{0}}\right\}
\end{equation}
the $n$-dimensional {\em Krylov subspace\/} generated by repeated 
applications of ${\CB H}$ on the initial residual ${\bmi  r}_{0}$. 
The $2$ classes of methods considered here are distinguished by the choice 
of the
variational functional $J(\cdot)$.




\section [High order Schwinger-Levine approximations]
{High order Schwinger-Levine approximations of the scattering 
transition using Pad\'{e} approximation theory} \label{sectionpade}

In this section, we show how to construct the high order approximations
to the {\em scattering transition\/}
\begin{equation} \label{scatTrans}
\langle \tilde{\bmi  b}, {\bmi  v}\rangle
\equiv
\langle \tilde{\bmi  b}, {\bmi  \Omega} {\bmi  b}\rangle 
\end{equation}
between a given incident state ${\bmi  b}$ and scattered state 
$\tilde{\bmi  b}$
using Pad\'{e} theory. The  M{\o}ller operator ${\bmi  \Omega}$ was given 
in 
(\ref{moller}), and the bilinear form $<\cdot,\cdot>$ in
(\ref{scatTrans}) was given  by (\ref{bilinearDef}). The presentation of 
this thesis was inspired by {\em Nuttall\/}~(1970). The material 
now appears to be well known within certain scientific disciplines 
[e.g. {\em Garibotti\/}~(1972); {\em Lucchese, Takatsuka, \& 
McKoy\/}~(1986); {\em Duneczky \& Wyatt\/}~(1988)].

For notational convenience we introduce the {\em generalized moment\/} 
${ M}_{  n}$ defined by
\begin{equation} 	\label{moment}
{  M}_{  n}=\left\langle\tilde{\bmi   b}\  , \
	{\left(\rf{\CB P}_{r,s} {\mit \Delta}{\CB C}_{eff}\right)}^{n}\ 
	{\bmi  b}\right\rangle
\end{equation}
Then using (\ref{moment}), the  leading order approximation
\begin{equation} \label{pade01}
	{[0,1]}_{<\tilde{\bmi  b},{\bmi  
v}>}={\frac{{M}_{0}^{2}}{{M}_{0}-{M}_{1}}}
\end{equation}
to (\ref{scatTrans}) was 
introduced into computational seismology by {\em Lapwood \& Hudson\/} 
~(1975).
We extend their method handle to heterogeneous media, and show how to 
compute the higher order corrections needed to secure a reasonable level 
of convergence to the true transition.

We begin with the simplest case of {\em weak scattering\/}. The scattering 
strength can be quantified by the {\em spectral radius\/} $\rho (\rf{\CB 
P}_{r,s}  
{\mit \Delta}{\CB C}_{eff} )$ 
of the {\em rescattering operator\/} $\rf{\CB P}_{r,s} {\mit \Delta}{\CB 
C}_{eff} $ 
(MacMillan \& Redish, 1986), defined as the maximum 
complex plane modulus of any eigenvalue 
$\lambda(\rf{\CB P}_{r,s}{\mit \Delta}{\CB C}_{eff} )$, of $\rf{\CB 
P}_{r,s}{\mit \Delta}{\CB C}_{eff}$ i.e.
\begin{equation} \label{rhoPS}
	\rho (\rf{\CB P}_{r,s}  {\mit \Delta}{\CB C}_{eff} )=\mbox{max} 
|\lambda(\rf{\CB P}_{r,s}  
	{\mit \Delta}{\CB C}_{eff})| .
\end{equation}
For the very pleasant case of weak scattering
\begin{equation}	\label{weakscatter}
\rho (\rf{\CB P}_{r,s}  {\mit \Delta}{\CB C}_{eff} )	\ll 1 .
\end{equation}
More generally, if $\rho (\rf{\CB P}_{r,s}  {\mit \Delta}{\CB C}_{eff} )	< 
1$, then the 
M{\o}ller operator of (\ref{moller}) has the readily 
computable {\em multiple scattering series representation\/}
\begin{equation}	\label{MSmoller}
{\bmi  \Omega}=\sum\limits_{  n\ =\ 0}^{ \infty }
{\left({\rf{\CB P}_{r,s} {\mit \Delta}{\CB C}_{eff}}\right)}^{n}
\end{equation}
so that the $\left\langle\tilde{\bmi  b},{\bmi  v}\right\rangle$ 
transition of 
(\ref{scatTrans}) may be written 
\begin{eqnarray} \label{MStransition}
\left\langle\tilde{\bmi  b},{\bmi  v}\right\rangle & = &
\left\langle\tilde{\bmi  b},{\bmi  \Omega}{\bmi  b}\right\rangle\nonumber\\
&=&
\sum\limits_{  n\ =\ 0}^{ \infty }
\left\langle{\tilde{\bmi   b}\  , \
{\left({\rf{\CB P}_{r,s} {\mit \Delta}{\CB C}_{eff}}\right)}^{n}\ 
{\bmi  b}}\right\rangle\nonumber\\
&=&\sum\limits_{  n\ =\ 0}^{ \infty } {  M}_{  n}
\end{eqnarray}
with ${  M}_{  n}$ given in (\ref{moment}). When $\rho (\rf{\CB P}_{r,s}  
{\mit \Delta}{\CB C}_{eff} )	\ll 1$, the multiple scattering 
representation of 
(\ref{MStransition}) is the most effective means for the construction of 
$\left\langle\tilde{\bmi  b},{\bmi  v}\right\rangle$ as only a few terms 
need 
be 
summed in (\ref{MStransition}) to give a highly accurate approximation. 
Unfortunately, for realistic crustal scattering environments, where the 
scattering takes place on a length scale spanning several wavelengths of 
the incident field $\bmi  b$, the spectral radius $\rho (\rf{\CB P}_{r,s}  
{\mit \Delta}{\CB C}(k_y) )$ is often not small and (\ref{MStransition}) 
fails to converge.
To circumvent this apparent impasse, we propose to {\em analytically 
continue\/} (\ref{MStransition}) beyond its {\em domain of convergence\/}
 by resumming it as a sequence of {\em Pad\'{e} approximants\/}. An 
 excellent introduction to the Pad\'{e} resummation of divergent series is 
 given by {\em Bender \& Orszag  \/}~(1978), \S8.3.
%%%%%%%%%%%%%%%%%%%%%%%%%%%%%%%%%%%%%%%%%%%%%%%%%%%%%%%%%%%%%%%%%%%%%%%%%%%%%%%%%



To this end, we introduce a {\em continuation parameter \/} $\theta$ and 
consider the multiple scattering series of the hypothetical M{\o}ller 
operator
\begin{eqnarray}	\label{MSmoller1}
{\bmi  \Omega}_{\theta}&=&({\bmi  I}-\theta \rf{\CB P}_{r,s} {\mit 
\Delta}{\CB C}_{eff})^{-1}\nonumber\\
&=&\sum\limits_{  n\ =\ 0}^{ \infty }\theta^{n}{\left({\rf{\CB P}_{r,s} 
{\mit \Delta}{\CB C}_{eff}}\right)}^{n}
\end{eqnarray}
where $\theta$ is chosen so that $\rho (\theta \rf{\CB P}_{r,s} {\mit 
\Delta}{\CB C}_{eff} 
) < 1$. 
Upon using ${\bmi  \Omega}_{\theta}$ of (\ref{MSmoller1}) in place 
of ${\bmi  \Omega}$ in (\ref{MStransition}) we get
\begin{eqnarray} \label{MStransition2}
\left\langle\tilde{\bmi  b},{\bmi  v}_{\theta}\right\rangle & = &
\left\langle\tilde{\bmi  b},{\bmi  \Omega}_{\theta}\cdot{\bmi  
b}\right\rangle\nonumber\\
&=&\sum\limits_{  n\ =\ 0}^{ \infty } \theta^n {  M}_{  n}.
\end{eqnarray}
We can then investigate the limit
\begin{equation}	\label{transitionlimit}
	\left\langle\tilde{\bmi  b},{\bmi  v}\right\rangle=\lim\limits_{\theta \ 
\rightarrow\ 
	1}\left\langle\tilde{\bmi  b},{\bmi  v}_{\theta}\right\rangle 
\end{equation}
as the required analytic continuation of the earlier result 
(\ref{MStransition}).
 
Following {\em Baker \& Graves-Morris\/}~(1981), Part 1,  we denote the  
$\left[{L/M}\right]$ rational Pad\'{e} approximation to an arbitrary 
analytic function  
\[
f(z)=\sum\limits_{ n\ =\ 0}^{ \infty } { c}_{ n}\ {z}^{n}
\]
as
\[
\left[{L/M}\right]\ =\ {\frac{{a}_{0}+{a}_{1}z+\cdots 
+{a}_{L}{z}^{L}}{{b}_{0}+{b}_{1}z+\cdots +{b}_{M}{z}^{M}}}
\] 
where the undetermined parameters $\left\{ {a}_{j}\right\} $ and 
$\left\{ {b}_{j}\right\} $ are chosen so that
\begin{equation} \label{powerseries}
 \sum\limits_{ n\ =\ 1}^{ \infty } { c}_{ n}{z}^{n}\ =\ 
 {\frac{{a}_{0}+{a}_{1}z+\cdots +{a}_{L}{z}^{L}}
 {{b}_{0}+{b}_{1}z+\cdots +{b}_{M}{z}^{M}}}\ +\ 
 O\left({{z}^{L+M+1}}\right).
\end{equation} 
Using a rather elegant result due to Nuttall ({\em vid. Baker \& 
Graves-Morris \/} Part.1~(1981)), we  construct the
$\left[{L/M}\right]$ approximation to (\ref{powerseries}), as an explicit 
function of the
 $\left\{ {c}_{j}\right\} $ according to
\begin{eqnarray} \label{nuttalform}
\lefteqn{\left[{L/M}\right]=  \sum\limits_{ n\ =\ 0}^{ L-M} 
{ c}_{ n} \ {z}^{n}\ +\ 
\left({\begin{array}{cccc}{c}_{L-M+1}&{c}_{L-M+2}&\cdots 
&{c}_{L}\end{array}}\right) \times} \nonumber\\
&&
\!\!\!\!\!\!\!\!\!\!\!\!\!\!
{\left({\begin{array}{cccc}{c}_{L-M+1}-z\ {c}_{L-M+2}&{c}_{L-M+2}-z\ 
{c}_{L-M+3}&\cdots &{c}_{L}-z\ {c}_{L+1}\\ {c}_{L-M+2}-z\ 
{c}_{L-M+3}&{c}_{L-M+3}-z\ {c}_{L-M+4}&\cdots &{c}_{L+1}-z\ {c}_{L+2}\\ 
\vdots&\vdots&&\vdots\\ {c}_{L}-z\ {c}_{L+1}&{c}_{L+1}-z\ {c}_{L+2}&\cdots 
&{c}_{L+M-1}-z\ 
{c}_{L+M}\end{array}}\right)}^{-1}\left({\begin{array}{c}{c}_{L-M+1}\\ 
{c}_{L-M+2}\\ \vdots\\ {c}_{L}\end{array}}\right)\nonumber\\
&& 
\end{eqnarray}

By setting $L=Q-1$ and $M=Q$ in (\ref{nuttalform}),  the 
$\left[{(Q-1)/Q}\right]$ 
Pad\'{e} approximate  of
 (\ref{MStransition2}), denoted $\left[{(Q-1)/Q}\right]_{<\tilde{\bmi  
 b},{\bmi  v}>}$,  may be written
\begin{eqnarray} \label{nuttalform1}
\lefteqn{{\left[{(Q-1)/Q}\right]}_{<\tilde{\bmi  b},{\bmi  v}>}\ =
\ \ \left({\begin{array}{cccc}{M}_{0}&{M}_{1}&\cdots &{M}_{Q-1}\end{array}}
\right)\times} \nonumber\\
&&
{\left({\begin{array}{cccc}{M}_{0}-\theta \ 
{M}_{1}&{M}_{1}- \theta \ {M}_{2}&\cdots &{M}_{Q-1}- \theta 
 \ {M}_{Q}\\ {M}_{1}- \theta  \ {M}_{2}&{M}_{2}- 
\theta  \ {M}_{3}&\cdots &{M}_{Q}- \theta  \ {M}_{Q+1}\\ 
\vdots&\vdots&&\vdots\\ {M}_{Q-1}- \theta  \ {M}_{Q}&{M}_{Q}- 
\theta  \ {M}_{Q+1}&\cdots &{M}_{2 Q-2}- \theta  \ 
{M}_{2 Q-1}\end{array}}\right)}^{-1}\left({\begin{array}{c}{M}_{0}\\ 
{M}_{1}\\ \vdots\\ {M}_{Q-1}\end{array}}\right)
\end{eqnarray}
where the modified moments ${M}_{n}$ are given in (\ref{moment}). The 
analytic continuation of (\ref{transitionlimit}) is now readily 
accomplished by substituting
$\theta=1$ in (\ref{nuttalform1}), for appropriate values of $Q$. 
The $Q=1$ case is given by (\ref{pade01}). For  $Q=2$ and $\theta=1$,  
(\ref{nuttalform1}) reduces to
\begin{eqnarray} \label{pade12}
{[1,2]}_{<\tilde{\bmi  b},{\bmi  v}>} 
& = & \left({\begin{array}{cc}{M}_{0}&{M}_{1}\end{array}}\right)
{\left({\begin{array}{cc}{M}_{0}-{M}_{1}&{M}_{1}-{M}_{2}\\
{M}_{1}-{M}_{2}&{M}_{2}-{M}_{3}\end{array}}\right)}^{-1}
\left({\begin{array}{c}{M}_{0}\\ {M}_{1}\end{array}}\right) \nonumber\\
& = &\frac{M_0 M_{1}^2 +  M_{1}^3 
-  M_{0}^2 M_{2} - 2 M_{0} M_1 M_2 +  M_{0}^2 M_3}{ M_{1}^2 -  M_0 M_2 -  
M_1 M_2 +  M_{2}^2 + M_0 M_3 -  M_1 M_3},
\end{eqnarray}
while for $Q=3$ and $\theta=1$, (\ref{nuttalform1}) becomes
\begin{eqnarray} \label{pade23}
	{[2,3]}_{<\tilde{\bmi  b},{\bmi  v}>} & = & 
\left({\begin{array}{ccc}{M}_{0}&{M}_{1}&{M}_{2}\end{array}}\right){\left({\begin{array}{ccc}{M}_{0}-{M}_{1}&{M}_{1}-{M}_{2}&{M}_{2}-{M}_{3}\\ 


{M}_{1}-{M}_{2}&{M}_{2}-{M}_{3}&{M}_{3}-{M}_{4}\\ 
{M}_{2}-{M}_{3}&{M}_{3}-{M}_{4}&{M}_{4}-{M}_{5}\end{array}}\right)}^{-1}\left({\begin{array}{c}{M}_{0}\\ 


{M}_{1}\\ {M}_{2}\end{array}}\right) \nonumber\\
	 & = & \frac{N}{D}
\end{eqnarray}
where:
\begin{eqnarray}
\lefteqn{N=- M_{0} M_{2}^3   -  M_{1} M_{2}^3 - 
   M_{2}^4 + 2 M_{0} M_{1} M_{2} M_{3} + 
  2 M_{1}^2 M_{2} M_{3} +  M_{0} M_{2}^2 M_{3} +}\nonumber\\
&&\!\!\!\!\!\!\!\!\!\!\!\!\!\!
 3 M_{1} M_{2}^2 M_{3} -  M_{0}^2 M_{3}^2 - 
  2 M_{0} M_{1} M_{3}^2 -  M_{1}^2 M_{3}^2 - 
  2 M_{0} M_{2} M_{3}^2 -  M_{0} M_{1}^2 M_{4} - \nonumber\\
&& \!\!\!\!\!\!\!\!\!\!\!\!\!\! 
   M_{1}^3 M_{4} +  M_{0}^2 M_{2} M_{4} - 
  2 M_{1}^2 M_{2} M_{4} + 2 M_{0} M_{2}^2 M_{4} + 
   M_{0}^2 M_{3} M_{4} + 2 M_{0} M_{1} M_{3} M_{4} -\nonumber\\
&&\!\!\!\!\!\!\!\!\!\!\!\!\!\!
 M_{0}^2 M_{4}^2 +  M_{0} M_{1}^2 M_{5} + 
   M_{1}^3 M_{5} -  M_{0}^2 M_{2} M_{5} - 
  2 M_{0} M_{1} M_{2} M_{5} +  M_{0}^2 M_{3} M_{5};
\end{eqnarray}
\begin{eqnarray}
\lefteqn{D=  - M_{2}^3 + 2 M_{1} M_{2} M_{3} +  M_{2}^2 M_{3} - 
   M_{0} M_{3}^2 -  M_{1} M_{3}^2 -  M_{2} M_{3}^2 + 
   M_{3}^3 -  M_{1}^2 M_{4} +   }\nonumber\\
&&\!\!\!\!\!\!\!\!\!\!\!\!\!\!
 M_{0} M_{2} M_{4} - M_{1} M_{2} M_{4} +  M_{2}^2 M_{4} + 
   M_{0} M_{3} M_{4} +  M_{1} M_{3} M_{4} - 
  2 M_{2} M_{3} M_{4} -  M_{0} M_{4}^2 + \nonumber\\
&&\!\!\!\!\!\!\!\!\!\!\!\!\!\!
 M_{1} M_{4}^2 +  M_{1}^2 M_{5} -  M_{0} M_{2} M_{5} - 
   M_{1} M_{2} M_{5} +  M_{2}^2 M_{5} + 
   M_{0} M_{3} M_{5} -  M_{1} M_{3} M_{5}.
\end{eqnarray}

The procedure may be extended to higher values of $Q$ in 
(\ref{nuttalform1}), but explicit 
formulae are too cumbersome to write down. An alternate resummation
of the Born series approximation to the scattering transition 
(\ref{MStransition})
 based on the theory of {\em continued fractions\/} has been given by 
{\em Znojil\/}~(1986). Finally, both Pad\'{e} and continued fraction theory
have matrix generalizations, allowing us to model multiple scattering 
transitions simultaneously.

 In \S\ref{sectionGBCG}, we 
explore a mathematically equivalent but computationally superior
algorithm to compute the higher order corrections to the scattering 
transition (\ref{scatTrans}), based on a modified {\em bi-conjugate 
gradient recursion\/}. The relationship between Pad\'{e} approximation 
theory
and the biconjugate gradient method is now well known [e.g. {\em 
Gutknecht\/}~(1990); {\em Gutknecht\/}~(1992)].




\section[A variational summation formula]
{A variational summation formula for the scattering transition
derived from the bi-conjugate gradient method} \label{sectionGBCG}

Given that the {\em incident \/} state ${\bmi  b}$, and {\em scattered \/}
state $\tilde{\bmi  b}$, are each governed by the {\em Lippmann-Schwinger} 
type 
equations
\begin{equation} \label{primaldual}
{\CB H} \cdot  {\bmi  v} = {\bmi  b},\quad \mbox{and} \quad
{\CB H}^{\dagger} \cdot  \tilde{\bmi  v} = \tilde{\bmi  b}
\end{equation}
where ${\CB H}$ was given by (\ref{Hdef}), and the dual operator ${\CB 
H}^{\dagger}$ was given in (\ref{Htdef}).
We propose the {\em trial solutions\/}
\begin{eqnarray} \label{trialsolutions}
{\bmi  v}_{n}&=& \sum\limits_{  j\ =\ 0}^{ n } {\alpha }_{ 
 j}{ \bmi   p}_{  j}\\
{\tilde{\bmi  v}}_{n}&=& \sum\limits_{  j\ =\ 0}^{ n } 
{\tilde{\alpha }}_{  j}{\tilde{ \bmi   p}}_{  j}.
\end{eqnarray}
The scattering transition $\langle{\tilde{\bmi  b}, {\bmi  
v}}\rangle=
\langle{\tilde{\bmi  b}, {\bmi  \Omega} \cdot {\bmi  
b}}\rangle
= \langle{{\bmi  \Omega} \cdot\tilde{\bmi  b},  {\bmi  
b}}\rangle = \langle{\tilde{\bmi  v}, {\bmi  
b}}\rangle$ may be approximated as
\begin{equation} \label{BCGscat}
\left\langle{\tilde{\bmi  b}, {\bmi  \Omega} \cdot {\bmi  b}}\right\rangle 
\approx 
{\frac{{\left\langle{\tilde{\bmi  b},{\bmi  
b}}\right\rangle}^{2}}{\left\langle{\tilde{\bmi  b},{\CB H} 
\cdot {\bmi   b}}\right\rangle}}\ +\
\sum\limits_{  n\ =\ 1}^{ \infty } 
{\frac{\left\langle{{\tilde{ \bmi  r}}_{  n-1}  
,{\bmi  r}_{n-1}}\right\rangle^2}{\left\langle{\tilde{ \bmi  p}{}_{  
n}  ,{\CB H} \cdot { \bmi  p}_{  n}}\right\rangle}}
\end{equation} 
where ${\bmi  \Omega}$ is the {\em M{\o}ller\/} operator of 
(\ref{moller}), and
the {\em residual fields\/} ${\bmi  r}_{n}$ and $\tilde{\bmi  r}{}_{n}$ in 
(\ref{BCGscat})
are given by
\begin{eqnarray} \label{residuals}
{{\bmi  r}}_{n}&=& {{ \bmi   b}}  -{\CB H} \cdot {{  \bmi  v}}_{  
n}\nonumber\\ 
{\tilde{\bmi  r}}_{n}&=& {\tilde{ \bmi   b}}  -{\CB H}^{\dagger} \cdot 
{\tilde{\bmi  
v}}_{  n}. 
\end{eqnarray}

The term 
$
{{\left\langle{\tilde{\bmi  b},{\bmi  
b}}\right\rangle}^{2}}{\left\langle{\tilde{\bmi  b},{\CB H} \cdot 
{\bmi   b}}\right\rangle}^{-1}
$
in (\ref{BCGscat}) was given previously in
(\ref{pade01}) and was originally proposed by {\em Lapwood \& Hudson\/} 
~(1975). We will now demonstrate how the {\em high order additive 
corrections\/} in (\ref{BCGscat}) to the true transition 
$\langle{\tilde{\bmi  b}, {\bmi  
v}}\rangle$ are calculated in the course of the {\em Generalized 
bi-conjugate gradient \/}(GBCG) scheme. We have used the qualifier
``Generalized'' in the applellation ``Generalized bi-conjugate gradient 
scheme'' to highlight the nonstandard nature of the implementation adopted 
here. The standard formulation of {\em 
Lanczos\/}(1952), is based around the bilinear form 
$\left({\bmi  a}, {\bmi  b}\right)\equiv {\bmi  a}^{T}\cdot {\bmi  b}$, 
for a pair 
of complex vectors ${\bmi  a}$ and  ${\bmi  b}$. In our formulation, we 
employ
$\left\langle{\bmi  a}, {\bmi  b}\right\rangle\equiv 
{\bmi  a}^{T}\cdot  {\mit \Delta}{\CB C}_{eff} {\bmi  b}$, of 
(\ref{bilinearDef}), where the weight matrix has the 
symmetry $ {\mit \Delta}{\CB C}_{eff}(k_y) = {\mit \Delta}{\CB  
C}^{T}(-k_y) $. This 
modified choice of bilinear form is implicit in the work of {\em Lapwood 
\& Hudson\/}~(1975), although  explicit matrix 
representations are quite different from this thesis. 

While the choice of the weighted bilinear form (\ref{bilinearDef}), 
adversely affects the stability
of traditional interpretations of the BCG method based on vector updates 
to the solution field $\bmi  v$, the generalized bilinear form enhances 
the 
convergence of the variational scheme presented here, for specific 
scattering transitions  $\left\langle{\tilde{\bmi  b}, {\bmi  \Omega} 
\cdot 
{\bmi  b}}\right\rangle$. This point is rather subtle, but quite central 
to the 
modeling philosophy adopted here. The traditional BCG method, based on 
vector 
updates for the solution field $\bmi  v$ ``breaks down'' when the term
$\left\langle{{\tilde{ \bmi  r}}_{  n-1}  
,{\bmi  r}_{n-1}}\right\rangle$ in (\ref{BCGscat}) becomes so small that 
it is 
considerably contaminated by roundoff error. However, in the variational
representation of the specific transition $\left\langle{\tilde{\bmi  b}, 
{\bmi  \Omega} \cdot 
{\bmi  b}}\right\rangle$ of (\ref{BCGscat}), small values of 
$\left\langle{{\tilde{ \bmi  r}}_{  n-1}  
,{\bmi  r}_{n-1}}\right\rangle$ drive the correction scheme to convergence.
The important point that we wish to make at this point is:``although the 
scheme of this section bears a superficial resemblance to the standard 
BCG method, its numerical performance in realistic scattering experiments 
is substantially superior to the standard method, at the cost of 
concentrating on {\em one transition at a time\/}''.


To begin the derivation of (\ref{BCGscat}),
we introduce the {\em error functional\/}
\begin{equation} \label{errfunc}
	{\cal E}\left(\tilde{\bmi  v}{}_{n};{\bmi  
	v}{}_{n}\right)=\left\langle{\tilde{\bmi  v}-\tilde{\bmi  v}_{n}, 
	{\CB H} \cdot \left( {\bmi  v}-{\bmi  v}_{n} \right)}\right\rangle 
\end{equation}
where the trial solutions $\tilde{\bmi  v}{}_{n}$ and ${\bmi  v}_{n}$
were given in (\ref{trialsolutions}). We note that the functional 
(\ref{errfunc})
vanishes in the limit as ${\bmi  v}_n \rightarrow {\bmi  v}$ and
$\tilde{\bmi  v}_n \rightarrow \tilde{\bmi  v}$, where ${\bmi  v}$ and 
$\tilde{\bmi  v}$
are the exact solutions of (\ref{primaldual}). However, (\ref{errfunc}) 
does 
not define a norm since ${\cal E}\left(\tilde{\bmi  v}{}_{n};{\bmi  
	v}{}_{n}\right)=0$ does not imply  ${\bmi  v}_n={\bmi  v}$ and
$\tilde{\bmi  v}_n=\tilde{\bmi  v}$.
If we expand out the right hand side of (\ref{errfunc}) and rearrange the 
result, we readily 
find that
\begin{equation} \label{Jdef3}
	\left\langle{\tilde{\bmi  b}, {\bmi  \Omega} \cdot 
	{\bmi  b}}\right\rangle \ =\ {\cal E}\left(\tilde{\bmi  v}{}_{n};{\bmi  
	v}{}_{n}\right) \ +\  \left\langle\tilde{\bmi  b},{\bmi  
v}_{n}\right\rangle
\ +\   \left\langle\tilde{\bmi  v}{}_{n},{\bmi  b}\right\rangle
\ - \ \left\langle\tilde{\bmi  v}{}_{n},{\CB H} \cdot
{\bmi  v}_{n}\right\rangle .
\end{equation}
where the left hand side  term $\left\langle{\tilde{\bmi  v}, {\CB H} 
\cdot 
	{\bmi  v}}\right\rangle=
	\left\langle{\tilde{\bmi  b}, {\bmi  \Omega} 
\cdot 
	{\bmi  b}}\right\rangle$ in (\ref{Jdef3})  is the required {\em 
scattering 
	transition}. 
When the error functional ${\cal E}\left(\tilde{\bmi  v}{}_{n};{\bmi  
	v}{}_{n}\right)$ in (\ref{Jdef3}) is negligible, we get
\begin{eqnarray} \label{Jdef4}
	\left\langle{\tilde{\bmi  b}, {\bmi  \Omega} \cdot 
	{\bmi  b}}\right\rangle &\approx &   \left\langle\tilde{\bmi  
v}{}_{n},{\bmi  
b}\right\rangle
\ + \ \left\langle\tilde{\bmi  b},{\bmi  v}{}_{n}\right\rangle
\ - \ \left\langle\tilde{\bmi  v}{}_{n},{\CB H} \cdot
{\bmi  v}_{n}\right\rangle\nonumber\\
&\approx& {\cal J}(\tilde{\bmi  v}_n,\tilde{\bmi  b};{\bmi  v}_n, {\bmi  
b})
\end{eqnarray}

The {\em variational functional\/} ${\cal J}(\tilde{\bmi  
v}_n,\tilde{\bmi  
b};{\bmi  v}_n, {\bmi  b})$ in (\ref{Jdef4}) obeys the {\em addition rule}
\begin{eqnarray} \label{Jadd}
\lefteqn{{\cal 
J}({\tilde{\bmi  v}}_{n-1}+{\tilde{\bmi  z}}_{n},\tilde{\bmi  b};{\bmi  
v}_{n-1}+{\bmi  z}_{n},{\bmi  b})}\nonumber\\
&&={\cal J}({\tilde{\bmi  v}}_{n-1},\tilde{\bmi  b};{\bmi  v}_{n-1},{\bmi  
b})+\left\langle{\tilde{\bmi  b}-
{\CB H}^{\dagger} \cdot {\tilde{\bmi   v}}_{  n-1}  
,{\bmi  z}_{n}}\right\rangle+
\left\langle{{\tilde{\bmi  z}}_{n},{\bmi  b}-{\CB H} \cdot {  
\bmi  v}_{  n-1}}\right\rangle-\left\langle{{\tilde{\bmi  z}}_{n},{\CB H} 
\cdot 
{\bmi   z}_{  n}}\right\rangle\nonumber\\
&&=
{\cal J}({\tilde{\bmi  v}}_{n-1},\tilde{\bmi  b};{\bmi  v}_{n-1},{\bmi  
b})+
{\cal J}({\tilde{\bmi  z}}_{n},\tilde{\bmi  b}-{\CB H}^{\dagger}\cdot 
{\tilde{\bmi  
v}}_{n-1};
{{\bmi  z}}_{n},{\bmi  b}-{\CB H}\cdot {{\bmi  v}}_{n-1})
\end{eqnarray}
for arbitrary vectors ${\tilde{\bmi  z}}_{n}$ and ${{\bmi  z}}_{n}$. 

Assuming that the trial solutions $\tilde{\bmi  v}{}_{n}$ and ${\bmi  
v}{}_{n}$
in (\ref{Jdef4}) are given by (\ref{trialsolutions}), we can use the 
addition theorem
(\ref{Jadd}) to rewrite (\ref{Jdef4}) as
\begin{equation} \label{Jdef5}
	\left\langle{\tilde{\bmi  b}, {\bmi  \Omega} \cdot 
	{\bmi  b}}\right\rangle 
\approx \sum\limits_{ n\ =\ 0}^{\infty }
 {\cal J}({\tilde{\alpha}}_{n}{\tilde{\bmi  p}}_{n},{\tilde{\bmi  
r}}_{n-1};
{\alpha}_n{{\bmi  p}}_{n},{{\bmi  r}}_{n-1})
\end{equation}
where ${\tilde{\bmi  r}}_{n-1}=\tilde{\bmi  b}-{\CB H}^{\dagger}\cdot 
{\tilde{\bmi  
v}}_{n-1}$,
${{\bmi  r}}_{n-1}={\bmi  b}-{\CB H}\cdot {{\bmi  v}}_{n-1}$
with ${\tilde{\bmi  v}}_{-1}={{\bmi  v}}_{-1}={\bf  0}$ in (\ref{Jdef5}).

We now choose the {\em variational parameters\/} ${\tilde{\alpha}}_{j}$ and
${{\alpha}}_{j}$ in the trial solutions $\tilde{\bmi  v}{}_{n}$ and 
${\bmi  
v}{}_{n}$
so that
\begin{equation} \label{Jstat}
\sum\limits_{ n\ =\ 0}^{\infty }
 \delta{\cal J}({\tilde{\alpha}}_{n}{\tilde{\bmi  p}}_{n},{\tilde{\bmi  
r}}_{n-1};
{\alpha}_n{{\bmi  p}}_{n},{{\bmi  r}}_{n-1})=0
\end{equation}
which defines a {\em coupled\/} system for the variational parameters. 
A {\em sufficient condition\/} for (\ref{Jstat}) to hold is:
\[
\delta{\cal J}({\tilde{\alpha}}_{n}{\tilde{\bmi  p}}_{n},{\tilde{\bmi  
r}}_{n-1};
{\alpha}_n{{\bmi  p}}_{n},{{\bmi  r}}_{n-1})=0, \qquad n=0,1,\cdots
\]
so that the variational parameters may be 
determined recursively, i.e. 
${{\alpha}}_{j}$ defined in terms of 
$\{ {{\alpha}}_{0},\cdots,{{\alpha}}_{j-1} \}$ and similarly for 
${\tilde{\alpha}}_{j}$.

To begin, at recursion$-0$ we employ the {\em Born \/} approximation 
\begin{equation}
\tilde{\bmi  v}_{0}\equiv \tilde{\alpha}_{0}\tilde{\bmi  p}_{0}= 
\tilde{\alpha}_{0}\tilde{\bmi  b},
\quad \mbox{and} \quad
{\bmi  v}_{0}\equiv {\alpha}_{0}{\bmi  p}_{0}= {\alpha}_{0}{\bmi  b}
\end{equation}
with the variational parameters $\tilde{\alpha}_{0}$ and ${\alpha}_{0}$ 
chosen so that 
\begin{equation} \label{dJ0}
 \delta{\cal J}({\tilde{\alpha}}_{0}{\tilde{\bmi  p}}_{0},{\tilde{\bmi  
r}}_{-1};
{\alpha}_{0}{{\bmi  p}}_{0},{{\bmi  r}}_{-1})=\delta{\cal 
J}({\tilde{\alpha}}_{0}{\tilde{\bmi  b}},{\tilde{\bmi  b}};
{\alpha}_{0}{{\bmi  b}},{{\bmi  b}})=0
\end{equation}
and
\begin{equation} \label{J0}
{\cal 
J}({\tilde{\alpha}}_{0}{\tilde{\bmi  b}},{\tilde{\bmi  b}};
{\alpha}_0{{\bmi  b}},{{\bmi  
b}})=\tilde{\alpha}_{0}\left\langle\tilde{\bmi  
b},{\bmi  b}\right\rangle
\ +\  {\alpha}_{0} \left\langle\tilde{\bmi  b},{\bmi  b}\right\rangle
\ - \ \tilde{\alpha}_{0} {\alpha}_{0} \left\langle\tilde{\bmi  b},{\CB H} 
\cdot
{\bmi  b}\right\rangle.
\end{equation} 
To ensure the {\em stationary condition\/} (\ref{dJ0}), we impose
\begin{eqnarray} \label{a0}
{\frac{ \partial   \,}{ \partial   \,{ \alpha 
}_{0}}}{\cal 
J}({\tilde{\alpha}}_{0}{\tilde{\bmi  b}},{\tilde{\bmi  b}};
{\alpha}_0{{\bmi  b}},{{\bmi  b}})=0 &  
\Rightarrow  & {\tilde{ \alpha 
}}_{0}={\frac{\left\langle{\tilde{\bmi  b},{\bmi  
b}}\right\rangle}{\left\langle{\tilde{\bmi  b},{\CB H}
\cdot   {\bmi  b}}\right\rangle}}\nonumber\\
{\frac{ \partial   \,}{ \partial   \,\tilde{ \alpha 
}_{0}}}{\cal 
J}({\tilde{\alpha}}_{0}{\tilde{\bmi  b}},{\tilde{\bmi  b}};
{\alpha}_0{{\bmi  b}},{{\bmi  b}})=0 &  
\Rightarrow   & {{ \alpha 
}}_{0}={\frac{\left\langle{\tilde{\bmi  b},{\bmi  
b}}\right\rangle}{\left\langle{\tilde{\bmi  b},{\CB H}
\cdot   {\bmi  b}}\right\rangle}}
\end{eqnarray} 
so that substituting (\ref{a0}) for ${\tilde{\alpha}}_{0}$ and 
${\alpha}_{0}$
in (\ref{J0}) we find
\begin{equation} \label{J01}
{\cal 
J}({\tilde{\alpha}}_{0}{\tilde{\bmi  b}},{\tilde{\bmi  b}};
{\alpha}_0{{\bmi  b}},{{\bmi  b}})={\frac{\left\langle{\tilde{\bmi  
b},{\bmi  
b}}\right\rangle^2}{\left\langle{\tilde{\bmi  b},{\CB H}
\cdot   {\bmi  b}}\right\rangle}},
\end{equation} 
which is the first term on the right hand side of (\ref{BCGscat}). 
Eqn.(\ref{J01}) is essentially eqn.$(8.7)$ of {\em Lapwood \& Hudson\/} 
~(1975). 
These authors employed a physical argument based directly on {\em seismic 
reciprocity}.

At recursion-$n$, we choose the variational parameters $\tilde{ \alpha 
}_{n}$ and
${ \alpha }_{n}$ in the trial solutions
\begin{eqnarray} \label{ntrialsolutions}
{\bmi  v}_{n}&=& \sum\limits_{  j\ =\ 0}^{ n } {\alpha }_{ 
 j}{ \bmi   p}_{  j}\\
 &=& {\bmi  v}_{n-1}+{ \alpha }_{n}{\bmi  p}_{n}\nonumber\\
{\tilde{\bmi  v}}_{n}&=& \sum\limits_{  j\ =\ 0}^{ n } 
{\tilde{\alpha }}_{  j}{\tilde{ \bmi   p}}_{  j}\nonumber\\
&=& \tilde{\bmi  v}{}_{n-1}+\tilde{ \alpha }{}_{n}\tilde{\bmi  p}{}_{n}
\end{eqnarray}
so that
\begin{equation} \label{dJn}
 \delta{\cal J}({\tilde{\alpha}}_{n}{\tilde{\bmi  p}}_{n},{\tilde{\bmi  
r}}_{n-1};
{\alpha}_{n}{{\bmi  p}}_{n},{{\bmi  r}}_{n-1})=0
\end{equation}
where the {\em residuals\/} ${\tilde{\bmi  r}}_{n-1}$ and ${{\bmi  
r}}_{n-1}$
were given by (\ref{residuals}), so that
\begin{samepage}
\begin{eqnarray} \label{Jn}
\lefteqn{{\cal 
J}({\tilde{\alpha}}_{n}{\tilde{\bmi  p}}_{n},{\tilde{\bmi  r}}_{n-1};
{\alpha}_{n}{{\bmi  p}}_{n},{{\bmi  r}}_{n-1})
=}\nonumber\\
&&\qquad\qquad \tilde{\alpha}_{n}\left\langle\tilde{\bmi  p}{}_{n},{\bmi  
r}_{n-1}\right\rangle
\ +\  {\alpha}_{n} \left\langle\tilde{\bmi  r}{}_{n-1},{\bmi  
p}_{n}\right\rangle
\ - \ \tilde{\alpha}_{n} {\alpha}_{n} \left\langle\tilde{\bmi  p}{}_{n},
{\CB H} \cdot
{\bmi  p}_{n}\right\rangle.
\end{eqnarray} 
\end{samepage}
To ensure the {\em stationary condition\/} (\ref{dJn}), we impose
\begin{eqnarray} \label{an}
{\frac{ \partial   \,}{ \partial   \,{ \alpha 
}_{n}}}{\cal 
J}({\tilde{\alpha}}_{n}{\tilde{\bmi  p}}_{n},{\tilde{\bmi  r}}_{n-1};
{\alpha}_{n}{{\bmi  p}}_{n},{{\bmi  r}}_{n-1})=0 &  
\Rightarrow   & {\tilde{ \alpha 
}}_{n}={\frac{\left\langle{\tilde{\bmi  r}{}_{n-1},{\bmi  
p}_{n} }\right\rangle}{\left\langle{\tilde{\bmi  p}{}_{n},{\CB H}
\cdot   {\bmi  p}{}_{n}}\right\rangle}}\nonumber\\
{\frac{ \partial   \,}{ \partial   \,\tilde{ \alpha 
}_{n}}}{\cal 
J}({\tilde{\alpha}}_{n}{\tilde{\bmi  p}}_{n},{\tilde{\bmi  r}}_{n-1};
{\alpha}_{n}{{\bmi  p}}_{n},{{\bmi  r}}_{n-1})=0 &  
\Rightarrow   & {{ \alpha 
}}_{n}={\frac{\left\langle{\tilde{\bmi  p}{}_{n},{\bmi  
r}_{n-1}}\right\rangle}{\left\langle{\tilde{\bmi  p}{}_{n},{\CB H}
\cdot   {\bmi  p}_{n}}\right\rangle}}
\end{eqnarray} 
so that substituting (\ref{an}) for ${\tilde{\alpha}}_{n}$ and 
${\alpha}_{n}$
in (\ref{Jn}) we find
\begin{equation} \label{Jn1}
{\cal 
J}({\tilde{\alpha}}_{n}{\tilde{\bmi  p}}_{n},{\tilde{\bmi  r}}_{n-1};
{\alpha}_{n}{{\bmi  p}}_{n},{{\bmi  
r}}_{n-1})={\frac{\left\langle{\tilde{\bmi  
r}{}_{n-1},{\bmi  
p}_{n}}\right\rangle  
\left\langle{\tilde{\bmi  
p}{}_{n},{\bmi  
r}_{n-1}}\right\rangle }{\left\langle{\tilde{\bmi  p}{}_{n},{\CB H}
\cdot   {\bmi  p}_{n}}\right\rangle}}.
\end{equation} 
Finally, substituting (\ref{J01}) for 
\[
{\cal J}({\tilde{\alpha}}_{0} {\tilde{\bmi  p}}_{0} , {\tilde{\bmi  
r}}_{-1} ;
{\alpha}_{0} {{\bmi  p}}_{0}, {{\bmi  r}}_{-1})
\]
and (\ref{Jn1}) for 
\[
{\cal J} ({\tilde{\alpha}}_{n} {\tilde{\bmi  p}}_{n} , {\tilde{\bmi  
r}}_{n-1} 
;
{\alpha}_{n} {{\bmi  p}}_{n} , {{\bmi  r}}_{n-1}) 
\]
in the 
variational functional (\ref{Jadd}) gives the following {\em stationary\/} 
approximation to the $\langle{\tilde{\bmi  b}, {\bmi  \Omega} \cdot {\bmi  
b}}\rangle$ 
transition 
\begin{equation} \label{BCGscat0}
\left\langle{\tilde{\bmi  b}, {\bmi  \Omega} \cdot {\bmi  b}}\right\rangle 
\approx 
{\frac{{\left\langle{\tilde{\bmi  b},{\bmi  
b}}\right\rangle}^{2}}{\left\langle{\tilde{\bmi  b},{\CB H} 
\cdot {\bmi   b}}\right\rangle}}\ +\
\sum\limits_{  n\ =\ 1}^{ \infty } 
{\frac{\left\langle{{\tilde{ \bmi  p}}_{  n}  
,{\bmi  r}_{n-1}}\right\rangle\left\langle{{\tilde{ \bmi  r}}_{  n-1} 
 ,{\bmi  p}_{n}}\right\rangle}{\left\langle{\tilde{ \bmi  p}{}_{  
n}  ,{\CB H} \cdot { \bmi  p}_{  n}}\right\rangle}}.
\end{equation} 
The quality of the approximation (\ref{BCGscat0})
is dependent on 
the choice of the basis vectors $\tilde{\bmi  p}{}_{j}$ and 
${\bmi  p}{}_{j}$ in (\ref{trialsolutions}). An obvious choice are the 
{\em current steepest descent directions\/} 
\[
{\bmi  p}=-\nabla_{{\bmi  v}} {\cal J}=-(\tilde{\bmi  b}-{\CB 
H}^{\dagger}\cdot\tilde{\bmi  v})=-\tilde{\bmi  r}
\]
and
\[
\tilde{\bmi  p}=-\nabla_{\tilde{\bmi  v}} {\cal J}=-({\bmi  b}-{\CB 
H}\cdot{\bmi  v})=-{\bmi  r}
\]
so that we might choose ${\bmi  p}_{n}=-\tilde{\bmi  r}{}_{n-1}$ and 
$\tilde{\bmi  p}{}_{n}=-{\bmi  r}{}_{n-1}$. While this choice is 
reasonable, it does not reproduce the results of \S\ref{sectionpade}. 

Instead, we adopt the {\em basis\/} vectors $\tilde{\bmi  p}{}_{n}$ and 
${\bmi  
p}{}_{n}$
in the trial solution (\ref{trialsolutions}) 
commonly employed with the 
{\em bi-conjugate gradient method\/}, {\em viz.\/}
\begin{equation} \label{pupdate}
\left.\begin{array}{c}
{\bmi  p}_{n}= {\bmi  r}_{n-1}+{ \beta }_{n-1}{\bmi  p}_{n-1}\\
\tilde{\bmi  p}{}_{n}= \tilde{\bmi  r}{}_{n-1}+\tilde{ \beta 
}{}_{n-1}\tilde{\bmi  p}{}_{n-1}
\end{array}\right\}, \qquad n=1,2,\cdots
\end{equation}
with $\tilde{\bmi  p}{}_{0}=\tilde{\bmi  b}$ and ${\bmi  p}{}_{0}={\bmi  
b}$, 
consistent with the {\em Born approximation\/}
while the  residuals $\tilde{\bmi  r}{}_{n}$ and ${\bmi  r}_{n}$ were 
given 
previously in (\ref{residuals}) and may be updated from one iteration to 
the next via
\begin{eqnarray} \label{residualupdate}
{{\bmi  r}}_{n}&=& {{ \bmi   b}}  -{\CB H} \cdot \left({\bmi  v}_{  n-1} + 
\alpha_n {\bmi  p}_{  n} \right)
= {{ \bmi   r}}_{n-1}  - \alpha{}_n {\CB H} \cdot  {\bmi  p}{}_{  n}
\nonumber\\ 
{\tilde{\bmi  r}}_{n}&=& {\tilde{ \bmi   b}}  -{\CB H}^{\dagger} \cdot 
\left(\tilde{\bmi  
v}{}_{  n-1} + 
{\alpha}{}_n \tilde{\bmi  p}{}_{  n} \right)
= {\tilde{ \bmi   r}}_{n-1}  - {\alpha}{}_n {\CB H}^{\dagger} \cdot  
\tilde{\bmi  
p}{}_{  n}.
\end{eqnarray}

%%%%%%%%%%%%%%%%%%%%%%%%%%%%%%%%%%%%%%%%%%%%%%%%%%%%%%%%%%%%%%%%%%%%%%%%%%%%%%%%%%%%%%


The parameters $\tilde{ \beta }{}_{n-1}$ and ${ \beta }_{n-1}$ in 
(\ref{pupdate}) are 
chosen to enforce the {\em bi-conjugacy condition \/}
\begin{equation}	\label{Hconj}
	\left\langle{\tilde{\bmi  p}{}_{m}\ ,\ {\CB H} \cdot {\bmi  p}_{ 
n}}\right\rangle\ =
	\ 0,\ \ m \ne n. 
\end{equation}
Exploiting the considerable 
insight of {\em Lanczos\/}~(1952),  the {\em bi-conjugacy\/} condition 
(\ref{Hconj}) can be 
 achieved by ensuring that $\left\langle{\tilde{\bmi  
p}{}_{n-1}\ ,
 \ {\CB H} \cdot {\bmi  p}_{ n}}\right\rangle=
 \left\langle{\tilde{\bmi  
p}{}_{n}\ ,
 \ {\CB H} \cdot {\bmi  p}_{ n-1}}\right\rangle= 0$ or using 
 (\ref{pupdate}), this can be written as 
 \begin{eqnarray} \label{betadef}
\left\langle{\tilde{\bmi  p}{}_{n-1}\ ,
 \ {\CB H} \cdot \left({\bmi  r}_{n-1}+{ \beta }_{n-1}{\bmi  p}_{n-1} 
\right)}\right\rangle = 0 & \Rightarrow &
{ \beta }_{n-1} =\ -\ {\frac{\left\langle{\tilde{\bmi  p}{}_{n-1}\ ,
\ {\CB H}  \cdot {\bmi   r}_{ n-1}}\right\rangle}
{\left\langle{\tilde{\bmi  p}{}_{n-1}\ ,\ {\CB H} \cdot 
{\bmi   p}_{ n-1}}\right\rangle}}\nonumber\\
\left\langle{ \tilde{\bmi  r}{}_{n-1}+\tilde{ \beta }{}_{n-1}
\tilde{\bmi  p}{}_{n-1} \ ,
 \ {\CB H} \cdot {\bmi  p}_{ n-1}}\right\rangle = 0 &\Rightarrow&
\tilde{ \beta }{}_{n-1} =\ -\ {\frac{\left\langle{\tilde{\bmi  r}{}_{n-1}\ 
,
\ {\CB H}  \cdot {\bmi   p}_{ n-1}}\right\rangle}
{\left\langle{\tilde{\bmi  p}{}_{n-1}\ ,\ {\CB H} \cdot 
{\bmi   p}_{ n-1}}\right\rangle}}
\end{eqnarray}
The basic {\em variational summation\/} of (\ref{BCGscat0}) is   
summarized 
in Alg.${\bmi 6}$.${\bmi 1}$. The partial sum ${\mit \Delta} {  S}_{Q} $ 
at 
iteration $Q$ in 
Alg.${\bmi 6}$.${\bmi 1}$ is 
identical to the $\left[{(Q-1)/Q}\right]_{<\tilde{\bmi  
 b},{\bmi  v}>}$ {\em Pad\'{e} approximant\/} of (\ref{nuttalform1}) with 
the 
 continuation parameter set at $\theta=1$.

%%%%%%%%%%%%%%%%%%%%%%%%%%%%%%%%%%%%%%%%%%%%%%%%%%%%%%%%%%%%%%%%%%%%%%%%%%%%%%%%%
 


\begin{figure}
\begin{center}
\fbox{
\begin{minipage} {140.0mm}
{\bf  Algorithm }${\bmi 6}$.${\bmi 1}${\sc The basic variational 
summation for }$<\tilde{\bmi  b},{\bmi  v}>$\newline
\vspace{4mm}
\begin{itemize}

    \item  ${\mit \Delta} { S}_{-1} = 0 $
    
	\item  ${\bmi  p}_{0} = {\bmi  r}_{-1}= {\bmi  b}$
	
	\item  $\tilde{\bmi  p}{}_{0} = \tilde{\bmi  r}{}_{-1}= \tilde{\bmi  b}$

	\item {\bmi  for} $n=0,1,\ldots, $ until convergence
	
	\begin{itemize}
		\item  ${\alpha }_{n}  \ \ \  =\ {\frac{\left\langle{\tilde{\bmi  
p}{}_{n}\ ,\ 
{\bmi   r}_{n-1}}\right\rangle}{\left\langle{\tilde{\bmi  p}{}_{n}\ ,\ 
{\CB H} 
\cdot 
{\bmi  p}{}_{n}}\right\rangle}}$

		\item  $\tilde{\alpha }{}_{n}  \ \ \  =\ 
{\frac{\left\langle{\tilde{\bmi  
r}{}_{n-1}\ ,\ 
{\bmi   p}_{n}}\right\rangle}{\left\langle{\tilde{\bmi  p}{}_{n}\ ,\ {\CB 
H} 
\cdot 
{\bmi  p}{}_{n}}\right\rangle}}$
		
		\item  ${\mit \Delta}{S}_{n}   \ =\ {\mit \Delta}{ S}_{n-1}+{ \alpha 
}_{n}\langle
		\tilde{\bmi  r}{}_{n-1},{\bmi  p}_{n}\rangle$
		
\item  ${\bmi  r}_{n} \ \ \  =\ {\bmi  r}_{n-1}-{ \alpha }_{n} {\CB H} 
\cdot{\bmi  p}_{n}$

\item  $\tilde{\bmi  r}{}_{n} \ \ \  =\ \tilde{\bmi  r}{}_{n-1}-\tilde{ 
\alpha }{}_{n} {\CB H}^{\dagger} \cdot\tilde{\bmi  p}{}_{n}$
		
		\item  ${ \beta }_{n} \ \ \ =\ - {\frac{\left\langle{{\bmi  p}_{n}\ ,\ 
{\CB H}  \cdot {\bmi   r}_{ n}}\right\rangle}{\left\langle{{\bmi  p}_{n}\ 
,\ 
{\CB H} \cdot {\bmi   p}_{ n}}\right\rangle}}$
		
		\item $\tilde{ \beta }{}_{n} \ \ \ =\ - 
		{\frac{\left\langle{\tilde{\bmi  r}{}_{n}\ ,\ {\CB H}  \cdot {\bmi   
p}_{ n}}\right\rangle}
{\left\langle{\tilde{\bmi  p}{}_{n}\ ,\ {\CB H} \cdot {\bmi   p}_{ 
n}}\right\rangle}}$

		\item  ${\bmi  p}_{n+1} =\ {\bmi  r}_{n}+{ \beta }_{n}{\bmi  p}_{n}$

		\item  $\tilde{\bmi  p}{}_{n+1} =\ \tilde{\bmi  r}{}_{n}+\tilde{ \beta 
		}{}_{n}\tilde{\bmi  p}{}_{n}$

	\end{itemize}
	\item {\bmi  end}
	
	\item $\langle{\tilde{\bmi  b}, {\bmi  \Omega} \cdot {\bmi  
b}}\rangle\approx 
{\mit \Delta} {S}_{n}$
\end{itemize}
\end{minipage}
}
\end{center}
\end{figure}


%%%%%%%%%%%%%%%%%%%%%%%%%%%%%%%%%%%%%%%%%%%%%%%%%%%%%%%%%%%%%%%%%%%%%%%%%%%%%

Next, we derive the more computationally convenient summation algorithm of 
(\ref{BCGscat}) based on 
alternative formulae for the parameters $\alpha_n$ and $\beta_n$, that 
are identical to those employed in the {\em bi-conjugate gradient method}.
Using the update formula (\ref{pupdate})
in (\ref{an}) we find that
\begin{eqnarray} \label{eq3}
\lefteqn{
{{ \alpha }}_{n}={\frac{\left\langle{\tilde{\bmi  p}{}_{n},{\bmi  
r}_{n-1}}\right\rangle}{\left\langle{\tilde{\bmi  p}{}_{n},{\CB H}
\cdot   {\bmi  p}_{n}}\right\rangle}}=
\frac{\left\langle{{\tilde{ \bmi  r}}_{ n-1}+\tilde{\beta}{}_{n-1} 
{\tilde{ \bmi  
p}}_{ n-1} ,{\bmi  r}_{n-1}}\right\rangle}{\left\langle{\tilde{\bmi  
p}{}_{n},{\CB H}
\cdot   {\bmi  p}_{n}}\right\rangle}
=
\frac{\left\langle{{\tilde{ \bmi  r}}_{ n-1},{\bmi  r}_{n-1}}\right\rangle 
+
\tilde{\beta}{}_{n-1} \left\langle  {\tilde{ \bmi  
p}}_{ n-1} ,{\bmi  r}_{n-1} \right\rangle}{\left\langle{\tilde{\bmi  
p}{}_{n},{\CB H}
\cdot   {\bmi  p}_{n}}\right\rangle}
}\nonumber\\
&& \!\!\!\!\!\!\!\!\!\!
{\tilde{ \alpha 
}}_{n}={\frac{\left\langle{\tilde{\bmi  r}{}_{n-1},{\bmi  
p}_{n} }\right\rangle}{\left\langle{\tilde{\bmi  p}{}_{n},{\CB H}
\cdot   {\bmi  p}{}_{n}}\right\rangle}}=
{\frac{\left\langle{{\tilde{ \bmi  r}}_{  n-1} 
 ,{{ \bmi  r}}_{ n-1}+{\beta}{}_{n-1} {{ \bmi  p}}_{ n-1}}\right\rangle}
 {\left\langle{\tilde{\bmi  p}{}_{n},{\CB H}^{\dagger}\cdot   {\bmi  
p}{}_{n}}\right\rangle}}
 =
{\frac{\left\langle{{\tilde{ \bmi  r}}_{  n-1} 
 ,{{ \bmi  r}}_{ n-1}}\right\rangle+{\beta}{}_{n-1}
 \left\langle{{\tilde{ \bmi  r}}_{  n-1} 
 , {{ \bmi  p}}_{ n-1}}\right\rangle}
 {\left\langle{\tilde{\bmi  p}{}_{n},{\CB H}\cdot   {\bmi  
p}{}_{n}}\right\rangle}}
\nonumber\\
&&
\end{eqnarray}
but,
\begin{eqnarray} \label{eq1}
\left\langle{{\tilde{ \bmi  p}}_{ n-1} ,{\bmi  r}_{n-1}}\right\rangle
&=& \left\langle{
{\tilde{ \bmi  p}}_{ n-1} ,{\bmi  r}_{n-2} -
{\alpha}{}_{n-1} {\CB H}\cdot {{ \bmi  p}}_{ n-1}
}\right\rangle\nonumber\\
&=& \left\langle{
{\tilde{ \bmi  p}}_{ n-1} ,{\bmi  r}_{n-2} -
{\frac{\left\langle{\tilde{\bmi  p}{}_{n-1},{\bmi  
r}_{n-2}}\right\rangle}{\left\langle{\tilde{\bmi  p}{}_{n-1},{\CB H}
\cdot   {\bmi  p}_{n-1}}\right\rangle}} {\CB H}\cdot {{ \bmi  p}}_{ n-1}
}\right\rangle\nonumber\\
&=&0
\end{eqnarray}
and similarly
\begin{eqnarray} \label{eq2}
\left\langle{{\tilde{ \bmi  r}}_{ n-1} ,{\bmi  p}_{n-1}}\right\rangle
&=& \left\langle{
\tilde{\bmi  r}{}_{n-2} -
\tilde{\alpha}{}_{n-1} {\CB H}^{\dagger}\cdot {\tilde{ \bmi  p}}_{ n-1}, 
{{ \bmi  p}}_{ 
n-1} 
}\right\rangle\nonumber\\
&=& \left\langle{
\tilde{\bmi  r}{}_{n-2} -
{\frac{\left\langle{\tilde{\bmi  r}{}_{n-2},{\bmi  
p}_{n-1} }\right\rangle}{\left\langle{\tilde{\bmi  p}{}_{n-1},{\CB H}
\cdot   {\bmi  p}{}_{n-1}}\right\rangle}}
{\CB H}^{\dagger}\cdot {\tilde{ \bmi  p}}_{ n-1},
{{ \bmi  p}}_{ n-1} ,
}\right\rangle\nonumber\\
&=&0
\end{eqnarray}
so that using (\ref{eq1}) and (\ref{eq2}) in (\ref{eq3}) we get
\begin{equation} \label{eq4}
{{ \alpha }}_{n}={\tilde{ \alpha }}_{n}=
{\frac{\left\langle{\tilde{\bmi  r}{}_{n-1},{\bmi  
r}_{n-1}}\right\rangle}{\left\langle{\tilde{\bmi  p}{}_{n},{\CB H}
\cdot   {\bmi  p}_{n}}\right\rangle}}
\end{equation}

Next we derive an alternative to (\ref{betadef}) for ${ \beta }_{n}$,
but first  we need to investigate
$\left\langle{ \tilde{\bmi  r}{}_{n-1}\ ,\ {\bmi   r}_{ n} }\right\rangle$.
Performing $\left\langle \tilde{\bmi  r}{}_{n-1}\ ,\ 
(\ref{residualupdate}) \right\rangle$, we have
\begin{equation} \label{eq5}
\left\langle{ \tilde{\bmi  r}{}_{n-1}\ ,\ {\bmi   r}_{ n} }\right\rangle=
\left\langle{ \tilde{\bmi  r}{}_{n-1}\ ,\ 
{{ \bmi   r}}_{n-1}  - \alpha{}_n {\CB H} \cdot  {\bmi  p}{}_{  n} 
}\right\rangle
=\left\langle{ \tilde{\bmi  r}{}_{n-1}\ ,\ {\bmi   r}_{ n-1} }\right\rangle
\left(1-
{\frac{\left\langle{\tilde{\bmi  r}{}_{n-1},{\CB H}\cdot{\bmi  
p}_{n}}\right\rangle}{\left\langle{\tilde{\bmi  p}{}_{n},{\CB H}
\cdot   {\bmi  p}_{n}}\right\rangle}}
\right)
\end{equation}
where we have used (\ref{eq4}) for ${{ \alpha }}_{n}$ in (\ref{eq5}), but
using (\ref{pupdate}), the $\left\langle{\tilde{\bmi  p}{}_{n},{\CB H}
\cdot   {\bmi  p}_{n}}\right\rangle$ term in the denominator of 
(\ref{eq5}) may be rewritten as
\begin{equation} \label{eq6}
\left\langle{\tilde{\bmi  p}{}_{n},{\CB H}\cdot   {\bmi  
p}_{n}}\right\rangle=
\left\langle{\tilde{\bmi  r}{}_{n-1} + 
\tilde{\beta}{}_{n-1} \tilde{\bmi  p}{}_{n-1},{\CB H}\cdot {\bmi  
p}_{n}}\right\rangle
=\left\langle{\tilde{\bmi  r}{}_{n-1},{\CB H}\cdot {\bmi  
p}_{n}}\right\rangle
\end{equation}
where we have used the {\em bi-conjugacy\/} condition 
$\left\langle{\tilde{\bmi  p}{}_{n-1},{\CB H}\cdot {\bmi  
p}_{n}}\right\rangle=0$ of (\ref{Hconj}). It follows from (\ref{eq6}), 
that (\ref{eq5}) reduces down to
\begin{equation} \label{eq7}
\left\langle{ \tilde{\bmi  r}{}_{n-1}\ ,\ {\bmi   r}_{ n} }\right\rangle=0
\end{equation}
Now, using (\ref{residualupdate}),  ${ \beta }_{n}$ from (\ref{betadef}) 
may be rewritten as
\begin{equation} \label{betadef1}
{ \beta }_{n} =\ -\ {\frac{\left\langle{ \tilde{\bmi  p}{}_{n}\ ,
\ {\CB H}\cdot {\bmi   r}_{ n}}\right\rangle}
{\left\langle{\tilde{\bmi  p}{}_{n}\ ,\ {\CB H} \cdot 
{\bmi   p}_{ n}}\right\rangle}}
\ = \ -\ {\frac{\left\langle{{\CB H}^{\dagger}\cdot \tilde{\bmi  p}{}_{n}\ 
,
\  {\bmi   r}_{ n}}\right\rangle}
{\left\langle{{\CB H}^{\dagger}\cdot\tilde{\bmi  p}{}_{n}\ ,\  
{\bmi   p}_{ n}}\right\rangle}}
\ = \ -\ {\frac{\left\langle{{\tilde{ \bmi   r}}_{n}-{\tilde{ \bmi   
r}}_{n-1}\ ,
\  {\bmi   r}_{ n}}\right\rangle}
{\left\langle{{\tilde{ \bmi   r}}_{n}-{\tilde{ \bmi   r}}_{n-1}\ ,\  
{\bmi   p}_{ n}}\right\rangle}}
\end{equation}
but from (\ref{eq7}) we know that $\left\langle{ \tilde{\bmi  r}{}_{n-1}\ 
,\ {\bmi   r}_{ n} }\right\rangle=0$, while in (\ref{eq2}), we found
$\left\langle{{\tilde{ \bmi  r}}_{ n} ,{\bmi  p}_{n}}\right\rangle=0$ 
so that  (\ref{betadef1}) may be written as
\begin{equation} \label{betadef2}
{ \beta }_{n} \ =  \ {\frac{\left\langle{{\tilde{ \bmi   r}}_{n}\ ,
\  {\bmi   r}_{ n}}\right\rangle}
{\left\langle{{\tilde{ \bmi   r}}_{n-1}\ ,\  
{\bmi   p}_{ n}}\right\rangle}}
\ = \
\ {\frac{\left\langle{{\tilde{ \bmi   r}}_{n}\ ,
\  {\bmi   r}_{ n}}\right\rangle}
{\left\langle{{\tilde{ \bmi   r}}_{n-1}\ ,\  
{\bmi   r}_{ n-1} + {\beta}_{n-1} {\bmi   p}_{ n-1}}\right\rangle}}
\ = \
\ {\frac{\left\langle{{\tilde{ \bmi   r}}_{n}\ ,
\  {\bmi   r}_{ n}}\right\rangle}
{\left\langle{{\tilde{ \bmi   r}}_{n-1}\ ,\  
{\bmi   r}_{ n-1}}\right\rangle}}
\end{equation}
where we have used (\ref{pupdate}) and (\ref{eq2}) again.
The arguments used to derive  (\ref{betadef2}), are closely duplicated in 
the derivation of an alternate formula for $\tilde{\beta}{}_{n}$ in 
(\ref{betadef}). To save space, the result is stated without 
proof as 
\begin{equation} \label{betadef3}
\tilde{ \beta }{}_{n} \ =  \
\ {\frac{\left\langle{{\tilde{ \bmi   r}}_{n}\ ,
\  {\bmi   r}_{ n}}\right\rangle}
{\left\langle{{\tilde{ \bmi   r}}_{n-1}\ ,\  
{\bmi   r}_{ n-1}}\right\rangle}} = { \beta }{}_{n}
\end{equation}

Using the alternative formulae (\ref{eq4}) and (\ref{betadef3}), in 
Alg.${\bmi 6}$.${\bmi 1}$,
we get the more efficient scheme of Alg.${\bmi 6}$.${\bmi 2}$, that 
effectively 
sums (\ref{BCGscat}).

%%%%%%%%%%%%%%%%%%%%%%%%%%%%%%%%%%%%%%%%%%%%%%%%%%%%%%%%%%%%%%%%%%%%%%%%%%%%%%%%%
 


\begin{figure}
\begin{center}
\fbox{
\begin{minipage} {140.0mm}
{\bf  Algorithm }${\bmi 6}$.${\bmi 2}${\sc The enhanced variational 
summation for }$<\tilde{\bmi  b},{\bmi  v}>$\newline
\vspace{4mm}
\begin{itemize}

    \item  ${\mit \Delta} {S}_{-1} = 0 $
    
	\item  ${\bmi  r}_{-1}= {\bmi  b}$
	
	\item  $\tilde{\bmi  r}{}_{-1}= \tilde{\bmi  b}$

	\item{\bmi  f}{\bmi  o}{\bmi  r} $n=0,1,\ldots,$ until convergence
	
	\begin{itemize}
	
	   \item $\rho_{n-1}=\langle{\tilde{\bmi  r}{}_{n-1}, {\bmi  
r}{}_{n-1}}\rangle$
	   \item{\bmi  i}{\bmi  f} $n=1$
	   \begin{itemize}
	   \item ${\bmi  p}{}_{n}={\bmi  r}{}_{n-1}$
	   \item $\tilde{\bmi  p}{}_{n}=\tilde{\bmi  r}{}_{n-1}$
	   \end{itemize}
	   \item{\bmi  e}{\bmi  l}{\bmi  s}{\bmi  e}
	   \begin{itemize}
	   \item $\beta_{n-1}=\rho_{n-1}/\rho_{n-2}$
	   \item ${\bmi  p}{}_{n}={\bmi  r}{}_{n-1}+\beta_{n-1} {\bmi  p}{}_{n-1}$
	   \item $\tilde{\bmi  p}{}_{n}=\tilde{\bmi  r}{}_{n-1}+\beta_{n-1} 
\tilde{\bmi  p}{}_{n-1}$
	   \end{itemize}
       \item{\bmi  e}{\bmi  n}{\bmi  d}{\bmi  i}{\bmi  f}
       \item ${\bmi  q}{}_{n}={\CB H} \cdot {\bmi  p}{}_{n}$       
       \item $\tilde{\bmi  q}{}_{n}={\CB H}^{\dagger} \cdot \tilde{\bmi  
p}{}_{n}$
       
	   \item  ${\alpha }_{n}  \ \ \  =\ {\frac{\rho_{n-1}}
	{\left\langle{\tilde{\bmi  p}{}_{n-1}\ ,\ {\bmi  q}{}_{n}}\right\rangle}}$

		
		\item  ${\mit \Delta}{S}_{n}    =\ {\mit \Delta}{S}_{n-1}+{ \alpha 
}_{n}{ \rho }_{n-1}$
		
        \item  ${\bmi  r}_{n} \ \ \  =\ {\bmi  r}_{n-1}-{ \alpha }_{n} 
{\bmi  q}_{n}$

        \item  $\tilde{\bmi  r}{}_{n} \ \ \  =\ \tilde{\bmi  
r}{}_{n-1}-\tilde{\alpha }{}_{n} \tilde{\bmi  q}{}_{n}$
		
	\end{itemize}
	\item{\bmi  e}{\bmi  n}{\bmi  d}
	
	\item $\langle{\tilde{\bmi  b}, {\bmi  \Omega} \cdot {\bmi  
b}}\rangle\approx 
{\mit \Delta} {S}_{n}$
\end{itemize}
\end{minipage}
}
%\label{alg2}
\end{center}
\end{figure}


%%%%%%%%%%%%%%%%%%%%%%%%%%%%%%%%%%%%%%%%%%%%%%%%%%%%%%%%%%%%%%%%%%%%%%%%%%%%%

\section{The generalized minimal residual method}
In \S\ref{sectionGBCG} we developed a variational scheme to approximate 
the 
scattering transition  $\left\langle{\tilde{\bmi  b}, 
{\bmi  \Omega} \cdot {\bmi  b}}\right\rangle$. We will see in 
\S\ref{sectioncomparison}, that the variational
method is highly competitive with its computational alternatives, but at 
the considerable cost of concentrating on one transition at a time. Also, 
the variational scheme lacks a mathematically rigorous 
convergence theory. This is due to the choice of the error functional
(\ref{errfunc}), which does not define a norm.


During the early stages of 
the development of the variational method  of \S\ref{sectionGBCG}, 
the author implemented a complex arithmetic version of the standard {\em 
Generalized minimal 
residual method\/} (GMRES) of {\em Saad \& Schultz}~(1986), to check the 
convergence of the variational method. The spinoff
from this development was a modified method that could compute all the 
desired scattering transitions at once, at an acceptable convergence rate.
Also, the method tends to be rather robust, even for  
scattering problems where the spectral radius $\rho(\rf{\CB P}_{r,s}{\mit 
\Delta}{\CB C}_{eff})> 1$, 
and the error control tends to reflect the true 
error, rather than some fanciful mathematical notion of that quantity. 
These 
qualities suggest that the GMRES method is a robust computational 
tool in its own right. The author is unaware of any published results on 
the application of the 
GMRES method to the solution of acoustic, electromagnetic or elastic wave 
scattering, with previous computational experience being confined to 
fluid dynamics and meteorology [e.g. {\em Kadio\u{g}lu \& Mudrick \/} 
~\shortcite{KadiogluandMudrick}]. In these published applications,  
the field variables are {\em real vectors\/}.

An elementary introduction to the  GMRES method is presented nicely  in 
({\em Bruaset\/}~ \shortcite{Bruaset}). An {\tt f77}  code for  real 
arithmetic 
problems, based directly the standard method proposed by {\em Saad \& 
Schultz}~(1986)
is available through {\tt Netlib}  [see also {\em Barrett et 
al.\/}~(1994)]. 
At the most basic level, the GMRES computes a sequence of vector updates 
to the approximate solution field, so that the correction introduced at 
each update minimizes the {\em least squares measure of the equation 
residual\/}. The attendant linear least squares problem is solved via QR 
factorizations. Each iteration of  basic GMRES method 
requires both an {\em outer\/} ({\em see\/} {\sc Box} {\bf 6.1}) and an 
{\em inner\/} ({\em see\/} {\sc Box} {\bf 6.2})  QR factorization. 
The outer QR factorization, {\em orthogonalizes\/} the trial space basis, 
so that the least squares problem simplifies to the canonical {\em upper 
Hessenberg least squares problem\/} ({\em see\/} {\sc Box} {\bf 6.2}). 
In {\em Saad \& Schultz}~(1986), the outer QR factorization was performed
 using a {\em modified Gramm-Schmidt \/} procedure [e.g. {\em 
Watkins}~(1992)]. 
 However, we have 
found that it is preferable to implement the outer QR factorization using 
a 
sequence of {\em unitary Householder reflections}, which are 
substantially more stable than modified Gramm-Schmidt factorizations,
especially for problems that require several iterations to secure 
convergence.
The second extension of the basic GMRES method concerns the choice of the 
{\em trial space basis\/}. In {\em Saad \& Schultz}~(1986), the standard 
{\em 
Krylov} trial space of (\ref{Krylovdef}) is used. In the outer QR 
factorization mentioned above, we  orthogonalize this basis set.
Following {\em Walker\/}~(1988), we select our trial space basis directly 
from the orthogonalized Krylov basis, rather than the original basis 
(\ref{Krylovdef}).  These two extensions of the basic method significantly
 enhance the observed convergence rate 
(\S\ref{sectioncomparison}).

Unlike the variational treatment of \S\ref{sectionGBCG}, the GMRES 
development does not introduce the {\em dual\/} problem
${\CB H}^{\dagger}\tilde{\bmi  v}=\tilde{\bmi  b}$, but is based solely 
on the {\em 
primal\/}  $N \times N$ matrix/vector equation 
\begin{equation} \label{primal}
{\CB H}{\bmi  v}={\bmi  b}. 
\end{equation}
where $N=8 N_x L$,  ${\CB H}$ was given by (\ref{Hdef}) and ${\bmi  b}$ 
was given by 
(\ref{bvardef})
Given the starting approximation ${\bmi  v}_{0}\in C^{N}$, to the solution 
of
(\ref{primal}), we propose to find the {\em correction\/} ${\bmi  z}_{n}$,
in the current approximate solution  ${\bmi  v}_{n}={\bmi  v}_{0}+{\bmi  
z}_{n}$ at iteration-$n$ so 
that the {\em residual\/} ${\bmi  r}_{n}\equiv{\bmi  b}-{\CB H}{\bmi  
v}_{n}$ is 
minimized over the $n$-dimensional {\em Krylov subspace\/} ${\cal 
K}_{n}({\bmi  r}_{0},
{\CB H})$
 of (\ref {Krylovdef}). We recall that ${\cal K}_{n}({\bmi  r}_{0},{\CB 
H})$ was 
 generated by repeated applications of ${\CB H}$ on the initial residual 
${\bmi  r}_{0}\equiv {\bmi  b}-{\CB H} \cdot {\bmi  v}_{0}$. The {\em 
minimal 
residual update\/} ${\bmi  z}_{n}$, is the  extremal solution of the
{\em least squares functional for the equation residual\/}, defined as
\begin{equation} \label{leastsquares}
J({\bmi  z}_{n},{\bmi  b}) \equiv  \min_{{\bmi  z}_{n}\in 
{\cal K}_{n}({\bmi  r}_{0},{\CB H})} 
\left|\!\left|
{{\bmi  b}-{\CB H} \cdot \left({\bmi  v}_{0} +{\bmi  z}_{n} 
\right)}\right|\!\right|_{2}.
\end{equation}
Since  ${\bmi  z}_{n}\in {\cal K}_{n}({\bmi  r}_{0},{\CB H})$, we also have
${\bmi  z}_{n}\in {\cal K}_{n}(\hat{\bmi  r}_{0},{\CB H})$, where 
where 
$\hat{\bmi  r}{}_{ 0}$ is the unit vector defined as  $\hat{\bmi  r}{}_{ 
0}= {\bmi  r}{}_{ 
0}/\| {\bmi  r}{}_{0} \|_{2}$, so that
\[
{\bmi  z}_{n}=\left[{\hat{\bmi  r}_{0},\ {\CB H} \cdot \hat{\bmi  r}_{ 0} 
,\cdots 
,{\CB H}^{n-1} \cdot \hat{\bmi  r}_{ 0}}\right]\cdot{\bmi  y}_{n}
\] 
where ${\bmi  y}_{n}$ is to be determined.
Hence, (\ref{leastsquares}) 
may be rewritten as
\begin{equation} \label{leastsquares2}
J({\bmi  y}_{n},{\bmi  b}) \equiv  \min_{{\bmi  y}_{n}\in C^{n}} 
\left|\!\left|
\left[{ \hat{\bmi  r}{}_{0},\ {\CB H} \cdot \hat{\bmi  r}{}_{ 0} ,\cdots 
,{\CB H}^{n} \cdot \hat{\bmi  r}{}_{ 0}}\right]\cdot\left(\begin{array}{c}
\| {\bmi  r}{}_{0} \|_{2}\\
-{\bmi  y}_{n}
\end{array} \right)
\right|\!\right|_{2}
\end{equation}




\begin{figure}
%\vspace*{10em}
{\em Orthonormalizing the trial space basis using unitary Householder 
reflections\/}
\fbox{
\begin{minipage} {140.0mm}
{\sc   Box} $\bf 6.1$ \newline

We will show how to construct the {\em Outer QR\/} factorization 
\[
\left[{ \hat{\bmi  r}{}_{0},\ {\CB H} \cdot \hat{\bmi  r}{}_{ 0} ,\cdots 
,{\CB H}^{n} \cdot \hat{\bmi  r}{}_{ 0}}\right] = 
{\bmi  P}_{1}{\bmi  P}_{2}\cdots {\bmi  P}_{n+1}\left[ \begin{array}{c} 
{\bmi R}_{n} \\ {\bf 0}_{(N-n-1) \times 
(n+1)}\end{array}  \right]
\]
to orthonormalize the  trial space basis ${\CB H}^{j} \cdot \hat{\bmi  
r}{}_{ 0}$,  
$N-$vectors.  The array ${\bmi  R}_{n}$ is an $(n+1) \times 
(n+1)$ upper triangular matrix, and ${\bmi  P}_{1}{\bmi  P}_{2} \cdots 
{\bmi  P}_{n+1}\equiv{\bmi  Q}_{n}$ is a factorization of the $N \times 
(n+1)$ array ${\bmi  Q}_{n}$ with orthonormal columns into the product of 
unitary Householder reflections. Also, in practice $N \gg (n+1)$.

$\quad$We begin with some basic results concerning unitary Householder 
reflections (Golub \& Van Loan , 1989; Bunch \& Grow, 1990; Watkins, 
1991). 
Let us assume that we have an $N-$vector ${\bmi m}=[m_1, m_2,
\cdots, m_N]^{T}$ with complex elements. Then the unitary Householder 
reflection
${\bmi P}_1$ is given by
\begin{equation} \label{P1def}
{\bmi P}_1\equiv {\bmi I}_N - \frac{2 {\bmi W}_{1} {\bmi W}_{1}^{H} 
}{{\bmi W}_{1}^{H}{\bmi W}_{1}}
\end{equation}
where 
\begin{equation} \label{W1def}
{\bmi W}_{1}=\left[ \begin{array}{c}m_1\\ m_2\\ \vdots\\ m_N \end{array} 
\right]+\begin{array}{cl}\left[ \begin{array}{c}\frac{\|[m_1, m_2,
\cdots, m_N]^{T}\|_{2} m_1}{|m_1|}\\ 0\\ \vdots\\ 0 \end{array} 
\right]&\begin{array}{c}1\\2\\ \vdots\\N \end{array} \end{array}
\end{equation}
so that ${\bmi P}_1$ transforms ${\bmi m}$ into
\[
{\bmi P}_1 {\bmi m}=-\left[ \begin{array}{c}\frac{\|[m_1, m_2,
\cdots, m_N]^{T}\|_{2} m_1}{|m_1|}\\ 0\\ \vdots\\ 0 \end{array}\right]
\]

Similarly, the unitary Householder reflection
${\bmi P}_2$ is given by
\begin{equation} \label{P2def}
{\bmi P}_2\equiv {\bmi I}_N - \frac{2 {\bmi W}_{2} {\bmi W}_{2}^{H} 
}{{\bmi W}_{2}^{H}{\bmi W}_{2}}
\end{equation}
where 
\begin{equation} \label{W2def}
{\bmi W}_{2}=\left[ \begin{array}{c}0\\ m_2\\ \vdots\\ m_N \end{array} 
\right]+\begin{array}{cl}\left[ \begin{array}{c}0\\
\frac{\|[0, m_2,
\cdots, m_N]^{T}\|_{2} m_2}{|m_2|}\\ 0\\ \vdots \end{array} 
\right]&\begin{array}{c}1\\2\\ \vdots\\N \end{array} \end{array}
\end{equation}
so that ${\bmi P}_2$ transforms ${\bmi m}$ into
\[
{\bmi P}_2 {\bmi m}=-\left[ \begin{array}{c}
m_1\\ \frac{\|[0, m_2,
\cdots, m_N]^{T}\|_{2} m_2}{|m_2|}\\ 0\\ \vdots \end{array}\right]
\]
and analogous definitions hold for ${\bmi P}_3, {\bmi P}_4, \cdots, {\bmi 
P}_{n+1}$.


\vspace{6 mm}
\end{minipage}
}
\end{figure}


\begin{figure}
%\vspace*{10em}
{\em Orthonormalizing the trial space basis using unitary Householder 
reflections\/}
\fbox{
\begin{minipage} {140.0mm}
{\sc   Box} $\bf 6.1$ {\em continued\/} \newline

To orthonormalize the trial space basis 
\[
\left[{ \hat{\bmi  r}{}_{0},\ {\CB H} \cdot \hat{\bmi  r}{}_{ 0} ,\cdots 
,{\CB H}^{n} \cdot \hat{\bmi  r}{}_{ 0}}\right]
\]
we operate with the Householder reflection ${\bmi P}_{1}$ to get
\begin{equation} \label{PK1}
\left[{ {\bmi P}_{1}\hat{\bmi  r}_{0},\ {\bmi P}_{1}{\CB H} \cdot 
\hat{\bmi  r}{}_{ 0} ,\cdots 
,{\bmi P}_{1}{\CB H}^{n} \cdot \hat{\bmi  r}{}_{ 0}}\right]
=
\left[{ -{\sigma}_{1} {\bmi e}_{1},\ {\bmi P}_{1}{\CB H} \cdot \hat{\bmi  
r}{}_{ 0} ,\cdots 
,{\bmi P}_{1}{\CB H}^{n} \cdot \hat{\bmi  r}{}_{ 0}}\right]
\end{equation}
where ${\sigma}_{1}= \| \hat{\bmi  r}_{0} \|_{2} (\hat{\bmi  r}_{0})_1/ 
|(\hat{\bmi  
r}_{0})_1| = (\hat{\bmi  r}_{0})_1/ |(\hat{\bmi  
r}_{0})_1|$, because $\| \hat{\bmi  r}_{0} \|_{2}=1$,  so that 
\[
-{\sigma}_{1} {\bmi e}_{1}= - \frac{ 
(\hat{\bmi  r}_{0})_1}{ |(\hat{\bmi  r}_{0})_1|} {\bmi e}_{1}
=-\left[ \begin{array}{c}\frac{ 
(\hat{\bmi  r}_{0})_1}{ |(\hat{\bmi  r}_{0})_1|}
\\ 0\\ \vdots\\ 0 \end{array}\right]
\]
and consequently the first coulumn of (\ref{PK1}) is in its final form.
We will now transform the second column of (\ref{PK1}) in such a way that 
the first column remains unaltered. To this end, we identify the second 
column in (\ref{PK1}) (i.e. ${\bmi P}_{1}{\CB H} \cdot \hat{\bmi  r}{}_{ 
0}$), with ${\bmi m}$ in (\ref{W2def}), so that we construct ${\bmi 
P}_{2}$ 
using its definition (\ref{P2def}) to get
\begin{equation} \label{P2def1}
{\bmi P}_2\equiv {\bmi I}_N - \frac{2 {\bmi W}_{2} {\bmi W}_{2}^{H} 
}{{\bmi W}_{2}^{H}{\bmi W}_{2}}
\end{equation}
where
\begin{equation} \label{W2def1}
{\bmi W}_{2}=\left[ \begin{array}{c}0\\ 
\left({\bmi P}_{1}{\CB H} \cdot \hat{\bmi  r}{}_{ 0} \right)_2\\ \vdots\\ 
\left({\bmi P}_{1}{\CB H} \cdot \hat{\bmi  r}{}_{ 0} \right)_N \end{array} 
\right]+\begin{array}{cl}\left[ \begin{array}{c}0\\
\frac{\|[0, \left({\bmi P}_{1}{\CB H} \cdot \hat{\bmi  r}{}_{ 0} \right)_2,
\cdots, \left({\bmi P}_{1}{\CB H} \cdot \hat{\bmi  r}{}_{ 0} 
\right)_N]^{T}\|_{2} 
\left({\bmi P}_{1}{\CB H} \cdot \hat{\bmi  r}{}_{ 0} \right)_2}
{|\left({\bmi P}_{1}{\CB H} \cdot \hat{\bmi  r}{}_{ 0} \right)_2|}\\ 0\\ 
\vdots \end{array} 
\right]&\begin{array}{c}1\\2\\ \vdots\\N \end{array} \end{array}
\end{equation}

Operating on (\ref{PK1}) with ${\bmi P}_2$ from (\ref{P2def1}), we get
\begin{eqnarray} \label{PK2}
\lefteqn{\!\!\!\!\!\!
\left[{ {\bmi P}_{1}\hat{\bmi  r}_{0},\ {\bmi P}_{2}{\bmi P}_{1}{\CB H} 
\cdot \hat{\bmi  r}{}_{ 0} ,\cdots 
,{\bmi P}_{2}{\bmi P}_{1}{\CB H}^{n} \cdot \hat{\bmi  r}{}_{ 0}}\right]=
\left[{ 
\left[\begin{array}{c}
-{\sigma}_{1}\\0\\0\\ 
\vdots \end{array}\right],\ 
\left[\begin{array}{c}
({\bmi P}_{1}{\CB H} \cdot \hat{\bmi  r}_{ 0})_{1}\\-{\sigma}_{2}\\0\\ 
\vdots \end{array}\right] ,\cdots 
,{\bmi P}_{2}{\bmi P}_{1}{\CB H}^{n} \cdot \hat{\bmi  r}{}_{ 
0}}\right]}\nonumber\\
&&
\end{eqnarray}
where 
$
{\sigma}_{2}= \frac{\| [0,({\bmi P}_{1}{\CB H} \cdot \hat{\bmi  r}_{ 
0})_{2}, \cdots, ({\bmi P}_{1}{\CB H} \cdot \hat{\bmi  r}_{ 
0})_{N}   ]^{T}  \|_{2} ({\bmi P}_{1}{\CB H} \cdot \hat{\bmi  r}_{ 
0})_{2}}{| ({\bmi P}_{1}{\CB H} \cdot \hat{\bmi  r}_{ 0})_{2}|}
$



$\quad$ Continuing in this fashion, we get
\begin{eqnarray} \label{PKn}
\lefteqn{\left[{ {\bmi P}_{1}\hat{\bmi  r}_{0},\ {\bmi P}_{2}{\bmi 
P}_{1}{\CB H} \cdot \hat{\bmi  r}{}_{ 0} ,\cdots 
,{\bmi P}_{2}{\bmi P}_{1}{\CB H}^{n} \cdot \hat{\bmi  r}{}_{ 0}}\right]=}\\
&&\!\!\!\!\!\!
\left[{ 
\left[\begin{array}{c}
-{\sigma}_{1}\\0\\0\\ \vdots\\ 0 \\ \vdots \\ 0 \end{array}\right],\ 
\left[\begin{array}{c}
({\bmi P}_{1}{\CB H} \cdot \hat{\bmi  r}_{ 0})_{1}\\-{\sigma}_{2}\\0\\ 
\vdots\\ 0 \\ \vdots \\ 0 \end{array}\right] ,\cdots, 
\left[\begin{array}{c}
({\bmi P}_{n}\cdots{\bmi P}_{1}{\CB H} \cdot \hat{\bmi  r}_{ 0})_{1}\\
({\bmi P}_{n}\cdots{\bmi P}_{2}{\CB H} \cdot \hat{\bmi  r}_{ 0})_{2}\\
\vdots\\ -{\sigma}_{n+1}\\ 0 \\ \vdots \\ 0 \end{array}\right] }\right]
\begin{array}{c}1 \\ 2 \\ \vdots \\ {n+1}\\ n+2 \\ \vdots \\ N \end{array}
= \left[ \begin{array}{c} {\bmi R}_{n} \\ {\bf 0}_{(N-n-1) \times 
(n+1)}\end{array}  \right] \nonumber
\end{eqnarray}

\vspace{2 mm}
\end{minipage}
}
\end{figure}

\begin{figure}
%\vspace*{10em}
{\em Orthonormalizing the trial space basis using unitary Householder 
reflections\/}
\fbox{
\begin{minipage} {140.0mm}
{\sc   Box} $\bf 6.1$ {\em continued\/} \newline

The decomposition (\ref{PKn}) may be written as
\begin{equation} \label{OuterQR0}
{\bmi  P}_{n+1}\cdots{\bmi  P}_{2}{\bmi  P}_{1}
\left[{ \hat{\bmi  r}{}_{0},\ {\CB H} \cdot \hat{\bmi  r}{}_{ 0} ,\cdots 
,{\CB H}^{n} \cdot \hat{\bmi  r}{}_{ 0}}\right] = 
\left[ \begin{array}{c} {\bmi R}_{n} \\ {\bf 0}_{(N-n-1) \times 
(n+1)}\end{array}  \right]
\end{equation}
but the unitary  Householder reflections ${\bmi  P}_{j}, 1 \le j \le 
(n+1)$) are {\em Hermitian\/} (i.e. ${\bmi  P}^{H}_{j}={\bmi  P}_{j}$), 
and 
consequently, they are {\em involutive\/} (i.e. ${\bmi  P}_{j}={\bmi  
P}^{-1}_{j}, 1 \le j \le (n+1)$), so we multiply both sides of 
(\ref{OuterQR0}) by 
\[
\left( {\bmi  P}_{n+1}\cdots{\bmi  P}_{2}{\bmi  P}_{1} \right)^{-1}=
{\bmi  P}_{1}{\bmi  P}_{2}\cdots{\bmi  P}_{n+1}
\]
to get
\begin{equation} \label{OuterQR1}
\left[{ \hat{\bmi  r}{}_{0},\ {\CB H} \cdot \hat{\bmi  r}{}_{ 0} ,\cdots 
,{\CB H}^{n} \cdot \hat{\bmi  r}{}_{ 0}}\right] = 
{\bmi  P}_{1}{\bmi  P}_{2}\cdots{\bmi  P}_{n+1}\left[ \begin{array}{c} 
{\bmi R}_{n} \\ {\bf 0}_{(N-n-1) \times 
(n+1)}\end{array}  \right]
\end{equation}
which is the required QR factorization of the trial space basis.

\vspace{2 mm}
\end{minipage}
}
\end{figure}


Using the {\em outer \/} QR factorization 
\[
\left[{ \hat{\bmi  r}{}_{0},\ {\CB H} \cdot \hat{\bmi  r}{}_{ 0} ,\cdots 
,{\CB H}^{n} \cdot \hat{\bmi  r}{}_{ 0}}\right] = {\bmi  Q}_{n}\left[ 
\begin{array}{c} {\bmi R}_{n} \\ {\bf 0}_{(N-n-1) \times 
(n+1)}\end{array}  \right]
\]
 ({\em see} {\sc Box} {\bf 6.1})  in
(\ref{leastsquares2}), we get 
\begin{eqnarray} \label{leastsquares3}
J({\bmi  y}_{n},{\bmi  b}) &\equiv&  \min_{{\bmi  y}_{n}\in C^{n}} 
\left|\!\left|
{\bmi  Q}_{n}\left[ \begin{array}{c} {\bmi R}_{n} \\ {\bf 0}_{(N-n-1) 
\times 
(n+1)}\end{array}  \right]\cdot\left(\begin{array}{c}
\| {\bmi  r}{}_{0} \|_{2}\\
-{\bmi  y}_{n}
\end{array} \right)
\right|\!\right|_{2}\nonumber\\
&\equiv & \min_{{\bmi  y}_{n}\in C^{n}} 
\left|\!\left|
\left[ \begin{array}{c} {\bmi R}_{n} \\ {\bf 0}_{(N-n-1) \times 
(n+1)}\end{array}  \right]\cdot\left(\begin{array}{c}
\| {\bmi  r}{}_{0} \|_{2}\\
-{\bmi  y}_{n}
\end{array} \right)
\right|\!\right|_{2}
\end{eqnarray}
where the second equation in (\ref{leastsquares3}) follows from the first 
equation, 
since the columns of  $N \times (n+1)$ array ${\bmi  Q}_{n}$ are {\em 
orthonormal\/} and consequently,  
$\left|\!\left|{\bmi  Q}_{n}\cdot {\bmi 
v}\right|\!\right|_{2}= \left|\!\left|{\bmi v}\right|\!\right|_{2}$ for 
any 
complex $(n+1)-$vector ${\bmi v}$. Since, 
\[
\left[ \begin{array}{c} {\bmi R}_{n} \\ {\bf 0}_{(N-n-1) \times 
(n+1)}\end{array}  \right]\equiv \left[-\frac{ 
(\hat{\bmi  r}_{0})_1}{ |(\hat{\bmi  r}_{0})_1|} {\bmi  e}_{1}, {\bmi 
H}_{n}\right]
\]
in (\ref{leastsquares3}) where ${\bmi H}_{n}$ is an
$(n+1)\times n$ upper Hessenberg array, the undetermined
vector ${\bmi  y}_{n}$ in (\ref{leastsquares3}) is the solution of the 
following {\em upper Hessenberg least squares problem}
\begin{equation} \label{leastsquares4}
J({\bmi  y}_{n},{\bmi  b}) \equiv  \min_{{\bmi  y}_{n}\in C^{n}} 
\left|\!\left|
 -\frac{ \| {\bmi  r}{}_{0} \|_{2}
(\hat{\bmi  r}_{0})_1}{ |(\hat{\bmi  r}_{0})_1|} {\bmi  e}_{1}-{\bmi 
H}_{n}\cdot{\bmi  y}_{n}
\right|\!\right|_{2}
\end{equation}
and the ${\bmi  e}_{1}$
is the $(n+1)$-vector ${\bmi  e}_{1}\equiv (1,0,\cdots,0)^{T}$.
To solve the linear least squares problem of (\ref{leastsquares4}) for
${\bmi  y}_{n}$, we decompose the   $(n+1)\times n$ Hessenberg array 
${\bmi  H}_{n}$
into its {\em inner\/} QR factors so that 
\[{\bmi  H}_{n}=\bar{\bmi  
Q}_{n}\left[\begin{array}{c} \bar{\bmi  R}_{n}\\ {\bf 0}_{1 \times n} 
\end{array}\right], 
\]
where ${\bf 0}_{1 \times n}$ is an extra row of zeros, 
$\bar{\bmi  Q}_{n}$ is a $(n+1)\times (n+1)$ array with orthonormal 
columns, and 
$\bar{\bmi  R}_{n}$
is an $n\times n$ upper triangular array. The {\em minimal residual 
solution \/}
${\bmi  y}_{n}$, of (\ref{leastsquares4}) , is now written as ({\em 
cf.\/}{\sc Box} {\bf 6.2})
\begin{equation} \label{minres} 
\left( \begin{array}{c}{\bmi  y}_{n}\\ 
0\end{array} \right)
=  -\frac{ \| {\bmi  r}{}_{0} \|_{2}
(\hat{\bmi  r}_{0})_1}{ |(\hat{\bmi  r}_{0})_1|} \ \left( 
\begin{array}{c}\bar{\bmi  R}_{n}^{-1}\\ 
{\bf 0}_{1 \times n}\end{array} \right)\bar{\bmi  Q}_{n}^{H}{\bmi  e}_{1}
\end{equation}
where  ${\bf 0}_{1 \times n}$   is the extra row of zeros, and the {\em 
equation 
residual\/} is given by
$\left|\!\left|{\bmi  r}\right|\!\right|_{2}\equiv  -\frac{ \| {\bmi  
r}{}_{0} \|_{2}
(\hat{\bmi  r}_{0})_1}{ |(\hat{\bmi  r}_{0})_1|} 
{\bmi  e}_{n+1}\cdot\bar{\bmi  Q}_{n}^{H}{\bmi  e}_{1}$. It is important 
to note 
that $\left|\!\left|{\bmi  r}\right|\!\right|_{2}$ is available without 
the 
need to construct  ${\bmi  R}^{-1}$ as in (\ref{minres}).

%%%%%%%%%%%%%%%%%%%%%%%%%%%%%%%%%%%%%%%%%%%%%%%%%%%%%%%%%%%%%%%%%%%%%%%%%%%%%

%%%%%%%%%%%%%%%%%%%%%%%%%%%%%%%%%%%%%%%%%%%%%%%%%%%%%%%%%%%%%%%%%%%%%%%%%%%%%

%%%%%%%%%%%%%%%%%%%%%%%%%%%%%%%%%%%%%%%%%%%%%%%%%%%%%%%%%%%%%%%%%%%%%%%%%%%%%

%%%%%%%%%%%%%%%%%%%%%%%%%%%%%%%%%%%%%%%%%%%%%%%%%%%%%%%%%%%%%%%%%%%%%%%%%%%%%


\begin{figure}
%\vspace*{10em}
{\em A  least squares solution of a full rank, complex, upper 
Hessenberg system using QR factors\/}
\fbox{
\begin{minipage} {140.0mm}
{\sc   Box} $\bf 6.2$ \newline


We begin with some elementary material concerning the application of {\em 
complex plane Given's 
rotations\/} [e.g. {\em Watkins\/}~(1991)] for the solution of the full 
rank, complex, upper 
Hessenberg least squares problem.

$\quad$ The $(n+1)\times n$ {\em upper-Hessenberg\/} array ${\bmi H}_{n}$ 
has the 
following sparsity pattern:
\[
{\bmi H}_{n}=\left[{\begin{array}{cccccc}
\ast&\ast&\ast&\cdots &\ast&\ast\\
\ast&\ast&\ast&\cdots &\ast&\ast\\
&\ast&\ast&\cdots &\ast&\ast\\
&&\ast&\cdots &\ast&\ast\\
&&&\ddots&\vdots&\vdots\\
&&&&\ast&\ast\\
&&&&&\ast\end{array}}\right]
\]
where each $\ast$ is a complex number.


$\quad$ Given a nonzero $\left[\begin{array}{c}a\\b\end{array}\right]\in 
C^{2}$ 
we define the {\em complex plane rotator\/} ${\bmi J}$ by
\[
{\bmi J}=\frac{1}{m }\left[\begin{array}{cc}
a&-{b}^{\ast}\\b&{a}^{\ast} \end{array}\right], \qquad m=\sqrt{|a|^2 + 
|b|^2}
\]
so that
\[
{\bmi J}^{H}\left[\begin{array}{c}a\\b\end{array}\right]=
\frac{1}{m }\left[\begin{array}{cc}
{a}^{\ast}&{b}^{\ast}\\-b&{a} 
\end{array}\right]\left[\begin{array}{c}a\\b\end{array}\right]
=\left[\begin{array}{c}m\\0\end{array}\right]
\]
Similarly, we define the complex rotator ${\bmi J}_{i+1,i}$ in the $x_i 
x_{i+1}$-plane so that
\begin{eqnarray} \label{Jijdef}
\lefteqn{{\bmi J}_{i+1,i}^{H}\left[{\begin{array}{c}{a}_{1,i}\\
\vdots\\
{a}_{i,i}\\
{a}_{i+1,i}\\
0\\
\vdots\\
0\end{array}}\right]=
\left[{\begin{array}{cccc}\begin{array}{cc}1&\\
&\ddots\end{array}&\begin{array}{cc}&\\
&\end{array}&\begin{array}{cc}&\\
&\end{array}&\begin{array}{c}\\
\end{array}\\
\begin{array}{cc}&\\
&\end{array}&{\frac{1}{{m}_{i+1}}}{\left[{\begin{array}{cc}{a}_{i,i}&-{a}_{i+1,i}^{\ast}\\

{a}_{i+1,i}&{a}_{i,i}^{\ast}\end{array}}\right]}^{H}&\begin{array}{cc}&\\
&\end{array}&\begin{array}{c}\\
\end{array}\\
\begin{array}{cc}&\\
&\end{array}&\begin{array}{cc}&\\
&\end{array}&\begin{array}{cc}1&\\
&\ddots\end{array}&\begin{array}{c}\\
\end{array}\\
\begin{array}{cc}&\end{array}&\begin{array}{cc}&\end{array}&
\begin{array}{cc}&\end{array}&1\end{array}}\right]
\left[{\begin{array}{c}{a}_{1,i}\\
\vdots\\
{a}_{i,i}\\
{a}_{i+1,i}\\
0\\
\vdots\\
0\end{array}}\right]=\left[{\begin{array}{c}{a}_{1,i}\\
\vdots\\
{m}_{i+1}\\
0\\
0\\
\vdots\\
0\end{array}}\right]}\nonumber\\
&&
\end{eqnarray}
which eliminates the $i+1$-element of the  vector 
\[
\left[{\begin{array}{ccccccc}{a}_{1,i},&
\cdots,&
{a}_{i,i},&
{a}_{i+1,i},&
0,&
\cdots,&
0\end{array}}\right]^{T}
\]
where $m_{i+1}=\sqrt{|{a}_{i,i}|^2 + |{a}_{i+1,i}|^2}$.


$\quad$ If we construct the $x_1 x_2$ rotator ${\bmi J}_{2,1}$ using the 
elements 
$({\bmi H}_{n})_{1,1}$ and $({\bmi H}_{n})_{2,1}$ in place of  ${a}_{i,i}$ 
and 
${a}_{i+1,i}$ respectively in (\ref{Jijdef}), and ${\bmi J}_{i+1,i}$ is 
constructed using the elements
\[
({\bmi J}^{H}_{i,i-1}\cdots {\bmi J}^{H}_{2,1}{\bmi H}_{n})_{i,i} \quad 
\mbox{and} \quad 
({\bmi J}^{H}_{i,i-1}\cdots {\bmi J}^{H}_{2,1}{\bmi H}_{n})_{i+1,i}
\]
in place of  ${a}_{i,i}$ and ${a}_{i+1,i}$ respectively in 
(\ref{Jijdef}), then
\[
\left[ \begin{array}{c}
\bar{\bmi R}_{n}\\{\bf 0}_{1 \times n}\end{array} \right]=
{\bmi J}^{H}_{n+1,n}{\bmi J}^{H}_{n,n-1}\cdots {\bmi J}^{H}_{2,1}{\bmi 
H}_{n}
\]

\vspace{6 mm}
\end{minipage}
}
\end{figure}



\begin{figure}
%\vspace*{10em}
\fbox{
\begin{minipage} {140.0mm}
{\sc   Box} $\bf 6.2$ {\em continued\/}\newline


or equivalently (because the rotators ${\bmi J}_{i,i-1}$ are {\em 
unitary\/}), we get the {\em QR factorization\/} of ${\bmi H}_{n}$ 
defined by
\[
{\bmi H}_{n} =
{\bmi J}_{2,1}{\bmi J}_{3,2}\cdots {\bmi J}_{n+1,n}
\left[ \begin{array}{c}\bar{\bmi R}_{n}\\{\bf 0}_{1 \times n}\end{array} 
\right]=
\bar{\bmi Q}_{n}\left[ \begin{array}{c}\bar{\bmi R}_{n}\\{\bf 0}_{1 \times 
n}\end{array} \right]
\]
where $\bar{\bmi Q}_{n}$ is the $(n+1)\times (n+1)$ unitary (so that 
${\bmi Q}^{H}_{n}{\bmi Q}_{n}={\bmi I}_{n+1}$) array
\[
\bar{\bmi Q}_{n}\equiv {\bmi J}_{2,1}{\bmi J}_{3,2}\cdots {\bmi J}_{n+1,n} 
\equiv
\left[\begin{array}{cccc} (\bar{\bmi Q}_{n})_{1},& 
(\bar{\bmi Q}_{n})_{2}& ,\cdots,& (\bar{\bmi Q}_{n})_{n+1}\end{array} 
\right]
\]
and $\bar{\bmi R}_{n}$ is the $n\times n$ {\em upper triangular matrix\/} 
with zeros 
below the diagonal, i.e.
\[
\left[ \begin{array}{c}\bar{\bmi R}_{n}\\{\bf 0}_{1 \times n}\end{array} 
\right]=
\left[{\begin{array}{cccccc}
\ast&\ast&\ast&\cdots &\ast&\ast\\
0&\ast&\ast&\cdots &\ast&\ast\\
0&0&\ast&\cdots &\ast&\ast\\
\vdots&\vdots&\vdots&\ddots&\vdots&\vdots\\
0&0&0&\cdots &\ast&\ast\\
0&0&0&\cdots &0&\ast\\
0&0&0&\cdots &0&0\end{array}}\right].
\]
The restriction that ${\bmi H}_{n}$ is 
of full column rank is required for the application of a QR factorization 
method.

$\quad$ Let us suppose that we are required to solve the {\em full column 
rank overdetermined 
linear system\/}
\[
{\bmi H}_{n}\cdot {\bmi y}_{n} = \bar{\bmi Q}_{n}
\left[ \begin{array}{c}\bar{\bmi R}_{n}\\{\bf 0}_{1 \times n}\end{array} 
\right]
\cdot {\bmi y}_{n}= {\bmi b}_{n}
\]
for the unknown $n-$vector ${\bmi y}_{n}$ given the $(n+1)-$element right 
hand side vector ${\bmi b}_{n}$. Using the unitarity of  $\bar{\bmi 
Q}_{n}$ we get
\[
\left[ \begin{array}{c}\bar{\bmi R}_{n}\\{\bf 0}_{1 \times n}\end{array} 
\right]\cdot 
{\bmi y}_{n}= \bar{\bmi Q}^{H}_{n}{\bmi b}_{n}= 
{\bmi J}^{H}_{n+1,n}{\bmi J}^{H}_{n,n-1}\cdots {\bmi J}^{H}_{2,1}{\bmi 
b}_{n}=
\left[\begin{array}{c}
(\bar{\bmi Q}^{H}_{n})_{1}\cdot{\bmi b}_{n}\\ 
(\bar{\bmi Q}^{H}_{n})_{2}\cdot{\bmi b}_{n}\\
\vdots\\
(\bar{\bmi Q}^{H}_{n})_{n+1}\cdot{\bmi b}_{n}
\end{array}  \right]
\]
and the corresponding {\em squared residual \/} is 
\[
\| \bar{\bmi Q}^{H}_{n}{\bmi b}_{n} - \left[ \begin{array}{c}\bar{\bmi 
R}_{n}\\
{\bf 0}_{1 \times n}\end{array} \right]\cdot 
{\bmi y}_{n} \|_{2}^{2}
\]



If we choose ${\bmi y}_{n}$ so that 
\[
{\bmi y}_{n}=
\bar{\bmi R}^{-1}_{n}\cdot 
\left[\begin{array}{c}
(\bar{\bmi Q}^{H}_{n})_{1}\cdot{\bmi b}_{n}\\ 
(\bar{\bmi Q}^{H}_{n})_{2}\cdot{\bmi b}_{n}\\
\vdots\\
(\bar{\bmi Q}^{H}_{n})_{n}\cdot{\bmi b}_{n}
\end{array}  \right]
\]
the corresponding squared residual 
$\| (\bar{\bmi Q}^{H}_{n})_{n+1}\cdot{\bmi b}_{n} \|_{2}^{2}$ is a minimum.





\vspace{6 mm}
\end{minipage}
}
\end{figure}

%%%%%%%%%%%%%%%%%%%%%%%%%%%%%%%%%%%%%%%%%%%%%%%%%%%%%%%%%%%%%%%%%%%%%%%%%%%%%

%%%%%%%%%%%%%%%%%%%%%%%%%%%%%%%%%%%%%%%%%%%%%%%%%%%%%%%%%%%%%%%%%%%%%%%%%%%%%

%%%%%%%%%%%%%%%%%%%%%%%%%%%%%%%%%%%%%%%%%%%%%%%%%%%%%%%%%%%%%%%%%%%%%%%%%%%%%

%%%%%%%%%%%%%%%%%%%%%%%%%%%%%%%%%%%%%%%%%%%%%%%%%%%%%%%%%%%%%%%%%%%%%%%%%%%%%




In the standard implementation of {\em Saad \& Schultz\/}~(1986), 
the 
\[
{\bmi  Q}_{n}\left[ \begin{array}{c} {\bmi R}_{n} \\ {\bf 0}_{(N-n-1) 
\times 
(n+1)}\end{array}  \right]
\]
 factorization is constructed using the {\em 
modified Gramm-Schmidt procedure\/}, while the 
\[
\bar{\bmi  Q}_{n}\left[ \begin{array}{c} \bar{\bmi R}_{n} \\ {\bf 0}_{1 
\times 
n}\end{array}  \right]
\]
 factorization
 is built from {\em unitary Givens rotations\/}.
 Unfortunately, the vectors in the  Krylov basis set (\ref{Krylovdef}) 
 become almost linearly dependent as the actual multiple scattering series 
 approaches convergence. The modified Gramm-Schmidt orthogonalization is 
 not robust enough to damp the roundoff inflicted by the increasingly 
 ill-conditioned Krylov basis set, and we have observed a dramatic 
improvement 
 in the numerical performance of the algorithm by replacing it with
 an orthogonalization based on {\em unitary Householder reflections\/} 
 instead. 





 
 Finally, we summarize the basic steps of the modified GMRES algorithm
 based on Householder reflections  in 
{\bf  algorithm }${\bmi 6}$.${\bmi 3}$. We have called the modified method
the GMRESH  algorithm, to distinguish it from the standard method.   



%%%%%%%%%%%%%%%%%%%%%%%%%%%%%%%%%%%%%%%%%%%%%%%%%%%%%%%%%%%%%%%%%%%%%%%%%%%%%

%%%%%%%%%%%%%%%%%%%%%%%%%%%%%%%%%%%%%%%%%%%%%%%%%%%%%%%%%%%%%%%%%%%%%%%%%%%%%

%%%%%%%%%%%%%%%%%%%%%%%%%%%%%%%%%%%%%%%%%%%%%%%%%%%%%%%%%%%%%%%%%%%%%%%%%%%%%

%%%%%%%%%%%%%%%%%%%%%%%%%%%%%%%%%%%%%%%%%%%%%%%%%%%%%%%%%%%%%%%%%%%%%%%%%%%%%

\begin{figure}
%\begin{center}
\[\!\!\!\!\!\!\!\!\!\!\!\!
\fbox{
\begin{minipage} {146.0mm}
{\bf  Algorithm }${\bmi 6}$.${\bmi 3}$ {\sc The  GMRESH algorithm}\newline

\begin{tt}
\noindent Suppose ${\bmi v}_{0}$, { Tol} and  { Maxit} are given.
\begin{enumerate}
\item  Initialize: Compute ${\bmi r}_{0}={\bmi b}-{\CB H}\cdot {\bmi 
v}_{0}$,
and determine  ${\bmi P}_{1}$ so that ${\bmi P}_{1}{\bmi r}_{0}=  -\frac{ 
\| {\bmi  r}{}_{0} \|_{2}
(\hat{\bmi  r}_{0})_1}{ |(\hat{\bmi  r}_{0})_1|} {\bmi e}_{1}=  \| {\bmi  
r}{}_{0} \|_{2} {\bmi v}_{1}$

\item Iterate: For n=1,Maxit, do:
\begin{enumerate}
\item Update Krylov basis: Set  ${\bmi v}_{n}={\bmi P}_{1}\cdots {\bmi 
P}_{n} {\bmi e}_{n}$
 (i.e. ${\bmi v}_{n}$ is the $n^{th}$ column of ${\bmi Q}_{n-1}$) and 
compute ${\CB H}\cdot {\bmi v}_{n}$

\item Outer QR factorization of Krylov basis: $[{\bmi v}_{1},{\CB H}\cdot 
{\bmi v}_{1},...,{\CB H}\cdot {\bmi v}_{n}]={\bmi Q}_{n} \left[ 
\begin{array}{c} {\bmi R}_{n} \\ {\bf 0}_{(N-n-1) \times 
(n+1)}\end{array}  \right]$, where
${\bmi R}_{n}=[  -\frac{ \| {\bmi  r}{}_{0} \|_{2}
(\hat{\bmi  r}_{0})_1}{ |(\hat{\bmi  r}_{0})_1|}{\bmi e}_{1},  {\bmi 
H}_{n}]$

\item Inner QR factorization: ${\bmi H}_{n}=\bar{\bmi Q}_{n} 
\left[ \begin{array}{c} \bar{\bmi R}_{n} \\ {\bf 0}_{1 \times 
n}\end{array}  \right]$

\item Find current residual: set  ${\bmi w}_{n}= -\frac{ \| {\bmi  
r}{}_{0} \|_{2}
(\hat{\bmi  r}_{0})_1}{ |(\hat{\bmi  r}_{0})_1|} \bar{\bmi Q}^{H}_{n} 
{\bmi e}_{1}$. Set current residual as $\mbox{res}_{n} =|{\bmi 
e}_{n+1}\cdot 
{\bmi w}_{n}|$.

\item \begin{enumerate}
\item If $\mbox{res}_{n}< {\tt Tol}$ then
\begin{enumerate}
\item go to {\bf 3}
\end{enumerate}

\end{enumerate}
\end{enumerate}


\item Update:
\begin{enumerate}
\item Set expansion coefficient vector ${\bmi  y}_{n}$ as
\[ 
\left( \begin{array}{c}{\bmi  y}_{n}\\ 
0\end{array} \right)
=  -\frac{ \| {\bmi  r}{}_{0} \|_{2}
(\hat{\bmi  r}_{0})_1}{ |(\hat{\bmi  r}_{0})_1|} \ \left( 
\begin{array}{c}\bar{\bmi  R}_{n}^{-1}\\ 
{\bf  0}_{1 \times n}\end{array} \right)\bar{\bmi  Q}_{n}^{H}{\bmi  e}_{1}
\]

\item ${\bmi v}_{0}\leftarrow {\bmi v}_{0}+[ {\bmi P}_{1}{\bmi e}_{1}, 
{\bmi P}_{1}{\bmi P}_{2}{\bmi e}_{2}, 
\cdots, {\bmi P}_{1}\cdots {\bmi P}_{n-1}{\bmi P}_{n}{\bmi e}_{n}]\cdot 
{\bmi  y}_{n} $


\item \begin{enumerate}
\item If $\mbox{res}_{n}< {\tt Tol}$ then
\begin{enumerate}
\item Accept ${\bmi v}_{0}$ as solution
\end{enumerate}
\item else
\begin{enumerate}
\item go to {\bf 1}

\end{enumerate}
\end{enumerate}


\end{enumerate}

\end{enumerate}
\end{tt} 
\vspace{2 mm} 
\end{minipage}
}
\]
%\end{center}
\end{figure}




\section[numerical comparison]
{A numerical comparison of the efficacy and efficiency  of the 
proposed schemes} \label{sectioncomparison}

We consider a stratified halfspace with a $14$km deep {\em monotone 
increasing velocity gradient zone\/} comprised of $80$ plane layers. We 
place the source and receiver in the surface layer. In each of the 
experiments, a {\em fast and dense\/} block $14 \times 14 \mbox{km}^2$ is 
superimposed on the reference medium with a given $\%$ perturbation with 
respect to the reference medium at each point. The dimensions of the block 
are chosen to span $10 \lambda_S \times 10 \lambda_S$ , where $\lambda_S$ 
is the effective SV-wavelength in the stratified reference medium averaged 
over the gradient zone at $2$Hz. We would expect to observe some very 
strongly scattered coherent phases in this experiment, because the 
heterogeneity is uniformly dense and fast. This problem is significantly 
more challenging to solve with an iterative scheme, than a similar sized 
random medium problem. In our problem, the scattered field from 
neighbouring 
regions has similar phase, leading to constructive interference. In 
random media, the phase from neighbouring points fluctuates, resulting in 
destructive interference, which essentially reduces the spectral radius 
of the problem, resulting in faster convergence of the multiple 
scattering series.  Results are only presented for the {\em in-plane P/SV 
problem\/} where $k_y=0$.

The numerical discretization  of the model problem has a matrix dimension 
of $N=57600$. To solve a 
problem of this size by storing the required $3.3 \times 10^{9}$ complex 
matrix 
elements (needing $\approx 53$ Gbytes) in an array and using some form of 
{\em direct method\/} based on 
{\em Gaussian elimination\/} is well beyond the capabilities of the 
present 
generation of super computers. 


The iterative schemes tested were:
\begin{itemize}

\item the {\em Born} series approximation of the scattering transition
(\ref{MStransition}). 

\item the analytically continued scattering transition of (\ref{BCGscat}),
computed via the  variational summation method of Alg.${\bmi 
6}$.${\bmi 2}$ 
(labeled {\sc BCGVAR\/} in the figures)

\item The basic {\sc GMRES\/} method of {\em Saad \& Schultz\/}~(1986).

\item The {\sc GMRESH} method

\item the standard {\em bi-conjugate gradient\/}  method of {\em 
Lanczos\/} 
~(1951) and {\em Lanczos\/} ~(1952)

\item the {\em stabilized bi-conjugate gradient} method of {\em van der 
Vorst\/}~(1992) [{\em see also\/} {\em Gutknecht\/}~ 
\shortcite{Gutknecht93}], that is essentially a hybrid version of the 
standard bi-conjugate gradient method and the GMRES(1) method (labeled 
{\sc 
BCGSTAB\/} in the figures). The standard algorithm has been adapted to 
handle complex arithmetic. 

\end{itemize} 

%%%%%%%%%%%%%%%%%%%%%%%%%%%%%%%%%%%%%%%%%%%%%%%%%%%%%%%%%%%%%%%%%%%%%%%%%%%%%%%%%%%%%%

\begin{figure}
\HideDisplacementBoxes
%\ShowDisplacementBoxes
\centerline{\BoxedEPSF{05.EPSF scaled 1000}}
\caption{Residual Vs. iteration for $5\%$ heterogeneity}
\protect\label{it05error}
\end{figure}
%%%%%%%%%%%%%%%%%%%%%%%%%%%%%%%%%%%%%%%%%%%%%%%%%%%%%%%%%%%%%%%%%%%%%%%%%%%%%


Some immediate conclusions follow from the experiment:
\begin{itemize}

\item the fastest methods for the {\em multi-transition\/} problem
are the {\sc GMRESH\/} and {\sc BCGSTAB\/} algorithms.

\item the variational summation was the fastest method 
( for a single transition), the $8$ iterations took approximately $18$ 
minutes on the DEC-$\alpha$.
This observation is consistent with calculations performed in {\em 
Wandzura \& Turley\/}~(1992), who considered biconjugate gradient 
based summations of the {\em electromagnetic mono-static scattering 
problem\/} (i.e. the analysis was restricted to $\tilde{\bmi b} = {\bmi 
b}$ transitions).  
We performed $8$ iterations for a relative error $<10^{-6}$ at $5\%$
perturbation, and $28$ iterations for a relative error $<10^{-6}$ at $10\%$
perturbation. In the $10\%$ problem, we are effectively performing an 
{\em analytic continuation\/} of the Born series transition beyond its 
domain of 
convergence, although this is at the cost of a substantial ``blow out"
in the required number of iterations.




\item the larger the contrast between the plane stratified 
reference medium and the actual heterogeneous medium (as measured by the 
$\%$ heterogeneity), the slower the 
convergence of the iterative schemes

\item the Born series converged the most slowly of the methods tested at 
$5\%$
heterogeneity and diverged at $10 \%$ heterogeneity. We observe that the 
first iterate of the Born series (usually referred to as the ``Born 
approximation"), is rather inadequate for the current problems. 



\item the standard bi-conjugate gradient method stalled, when approaching
the converged solution because the
 $\rho_{i}=<\tilde{\bmi  r}{}_{i},{\bmi  r}{}_{i}> \approx 10^{-14}$, was 
too small
and roundoff would have swamped to true solution.

\item the same parameter $\rho_i$ that caused the divergence in the 
bi-conjugate gradient method actually drives the variational summation 
method {\sc BCGVAR\/} to convergence, since the correction is 
$\frac{\rho^{2}_{i}}
{<\tilde{\bmi  p}{}_{i},{\CB H}\cdot{\bmi  p}{}_{i}>}$, and
$<\tilde{\bmi  p}{}_{i},{\CB H}\cdot{\bmi  p}{}_{i}>$ is not small in all 
the 
experiments conducted to date.
Indeed, small $\rho_{i}$ is a necessary, though not 
sufficient condition for the convergence of the variational method.

\item the convergence of  GMRES and GMRESH are {\em monotone}, with the 
later method markedly superior to the former in both tests.
Both  methods retain 
monotone convergence for more difficult problems as well, making 
them the method of choice for extremely nasty problems.

\end{itemize}

%%%%%%%%%%%%%%%%%%%%%%%%%%%%%%%%%%%%%%%%%%%%%%%%%%%%%%%%%%%%%%%%%%%%%%%%%%%%%%%%%%%%%%

\begin{figure}
\HideDisplacementBoxes
%\ShowDisplacementBoxes
\centerline{\BoxedEPSF{10.EPSF scaled 1000}}
\caption{Residual Vs. iteration for $10\%$ heterogeneity}
\protect\label{it10error}
\end{figure}
%%%%%%%%%%%%%%%%%%%%%%%%%%%%%%%%%%%%%%%%%%%%%%%%%%%%%%%%%%%%%%%%%%%%%%%%%%%%%


\section{The  preconditioned GMRESH method}
The most direct way to solve
\begin{equation} \label{primal1}
{\CB H}\cdot{\bmi  v}={\bmi  b}. 
\end{equation}
for the field vector ${\bmi  v}$, is to multiply both sides by the 
M{\o}ller operator 
\begin{equation} \label{Moll1}
{\bmi \Omega}={\CB H}^{-1}=\left( {\bmi I} - \rf{\CB P}_{r,s}\cdot {\mit 
\Delta}{\CB C}_{eff}   \right)^{-1}
\end{equation}
 to get
\begin{equation} \label{primal2}
{\bmi  v}={\bmi \Omega}\cdot{\bmi  b}. 
\end{equation}
but this method is much too expensive to be of practical use. However, 
it may be possible to construct an approximation $\tilde{\bmi \Omega}$ to 
${\bmi \Omega}$, and in instead of (\ref{Moll1}),  we solve the {\em 
preconditioned equation\/}
\begin{equation} \label{preconH}
\tilde{\bmi \Omega}{\CB H}\cdot{\bmi  v}=\tilde{\bmi \Omega}{\bmi  b}. 
\end{equation}
for the field vector ${\bmi  v}$. Provided that $\tilde{\bmi \Omega}$ is 
a reasonable approximation to ${\bmi \Omega}$, solving the preconditioned 
equation (\ref{preconH}) using the GMRES method will be more efficient 
in terms of the number of iterations required to secure convergence than 
solving 
the unpreconditioned equation (\ref{Moll1}).



To derive the preconditioning operator $\tilde{\bmi \Omega}$ we rewrite 
(\ref{Moll1}) as
\begin{equation} \label{Moll2}
{\bmi \Omega}=\left( {\bmi I} - \rf{\CB P}_{r,s}\cdot {\mit \Delta}{\CB 
C}_{eff}   \right)^{-1}
={\bmi I}+\rf{\CB P}_{r,s}\left( {\bmi I} - {\mit \Delta} {\CB C}\cdot  
\rf{\CB P}_{r,s}  \right)^{-1}{\mit \Delta} {\CB C}
={\bmi I}+\tilde{\CB P}{\mit \Delta} {\CB C}
\end{equation}
however $\tilde{\CB P}=\rf{\CB P}_{r,s}\big( {\bmi I} - {\mit \Delta} {\CB 
C}\cdot  \rf{\CB P}_{r,s}  
\big)^{-1}$ in (\ref{Moll2})
 is effectively a propagator for the 
actual medium because it models
all transmission and reflected phases in the 
 reference model, as well as the forward, backward and side scattered 
phases from the superimposed 
 heterogeneity. This operator could be constructed explicitly using 
matrix inversion  but this would be very 
expensive. Instead,  we will employ a forward scattering modification of
 the Riccati factorization algorithm 
 Alg.{\bf 5.3}  to construct the approximate propagator $\tilde{\CB P}_F$ 
that includes all 
 up/down transmission and reflected phases in the 
 reference model, and the forward scattered phases from the superimposed 
 heterogeneity, but {\em ignores backscattering and side scattering\/}.
   The preconditioning operator $\tilde{\bmi \Omega}$ is 
 then constructed explicitly as
\begin{equation} \label{Mollprecon}
\tilde{\bmi \Omega}=
{\bmi I}+\tilde{\CB P}_F {\mit \Delta} {\CB C}
\end{equation}
where $\tilde{\CB P}_F$ is constructed using Alg.{\bf 6.4}.

\begin{figure}
%\vspace*{10em}
\fbox{
\begin{minipage} {140.0mm} 
{\bf  Algorithm}~{\bf 6.4} {\sc A Riccati sweep implementation of the 
up/down  
preconditioning sweep \/} ${\CB B}=\tilde{\CB P}_F {\CB F}$\newline

We introduce the {\em forward scattering\/} operator ${\mit \delta}{\bmi 
S}_{\downarrow \downarrow }^{\rho}$  defined as
\begin{eqnarray} \label{laydSdefD}
\lefteqn{
{\mit \delta}{\bmi S}_{\downarrow \downarrow }^{\rho}=
-i\,\left({\begin{array}{cc}
{\bf 0}&\sqrt{{\bmi E}_{\rho}}
\end{array}}\right)\,
\hat{\bmi  D}{}^{T}_{z,\rho}(-{\bf k}_{x},-{k}_{y})}\\
&&
\times 
\left({\bmi I}_8\otimes{\bf DFT}^{T}(-{\bf k}_{x})  \right)
\hat{\bmi \Phi}{\mit \Delta}{\CB C}{}^{\rho}_{eff}(k_y) 
\hat{\bmi \Phi}{}^{-1} \left({\bmi I}_8\otimes 
{\bf DFT}({\bf k}_{x}) \right)\hat{\bmi  D}_{z,\rho}({\bf k}_{x},{k}_{y}) 
\left({\begin{array}{c}
\sqrt{{\bmi E}_{\rho}}\\ 
{\bf 0}
\end{array}}\right)\nonumber
\end{eqnarray}
and ${\mit \delta}{\bmi S}_{\uparrow \uparrow }^{\rho}$  defined as
\begin{eqnarray} \label{laydSdefU}
\lefteqn{
{\mit \delta}{\bmi S}_{\uparrow \uparrow }^{\rho}=
-i\,\left({\begin{array}{cc}
\sqrt{{\bmi E}_{\rho}}&{\bf 0}
\end{array}}\right)\,
\hat{\bmi  D}{}^{T}_{z,\rho}(-{\bf k}_{x},-{k}_{y})}\\
&&
\times 
\left({\bmi I}_8\otimes{\bf DFT}^{T}(-{\bf k}_{x})  \right)
\hat{\bmi \Phi}{\mit \Delta}{\CB C}{}^{\rho}_{eff}(k_y) 
\hat{\bmi \Phi}{}^{-1} \left({\bmi I}_8\otimes 
{\bf DFT}({\bf k}_{x}) \right)\hat{\bmi  D}_{z,\rho}({\bf k}_{x},{k}_{y}) 
\left({\begin{array}{c}
{\bf 0}\\ 
\sqrt{{\bmi E}_{\rho}}
\end{array}}\right)\nonumber
\end{eqnarray}
which are the {\em single scattering\/} analogues of 
 ${\mit \Delta}{\CB S}_{\downarrow \downarrow }^{\rho}(k_y)$ and  
 ${\mit \Delta}{\CB S}_{\uparrow \uparrow }^{\rho}(k_y)$
 respectively, from (\ref{laySdef}). Also, 
the {\em effective coupling matrix\/} ${\mit \Delta}{\CB 
C}{}^{\rho}_{eff}(k_y)$ 
for $\mbox{row}_{\rho}$ of the heterogeneity was given by (\ref{dCBCrho}).


\begin{itemize}
\item given the $8 L N_x$ {\em virtual source field\/} ${\CB F}$, compute 
the 
{\em preconditioned field\/} ${\CB B}=\tilde{\CB P}_F {\CB F}$


\begin{itemize}
\item  (*{\em bottom of layer stack\/ }*) 
\end{itemize}
\item ${\bmi R}_{ \uparrow \downarrow }^{L}\equiv 
{\bmi S}_{ \uparrow\downarrow }^{[L,\infty)}$

\item $
\left[{\begin{array}{c}{\bmi \Sigma}^{L}_{\uparrow}\\ 
{\bmi \Sigma}^{L}_{\downarrow}
\end{array}}\right]=-i\,\left({\begin{array}{cc}
\sqrt{{\bmi E}_{L}}&{\bf 0}\\ 
{\bf 0}&\sqrt{{\bmi E}_{L}}
\end{array}}\right)\,\hat{\bmi  D}{}^{T}_{z,L}(-{\bf 
k}_{x},-{k}_{y})\left({\bmi I}_8\otimes
{\bf DFT}^{T}(-{\bf k}_{x})  \right)\hat{\bmi \Phi}\left[{\CB F}_{L}  
\right]
$

\item  	$
{\bmi t}_{ \uparrow }^{L}=
{\bmi R}_{ \uparrow \downarrow }^{L}
{\bmi \Sigma}^{L}_{\downarrow}$

\begin{itemize}
\item  (*{\em upsweep\/ }*) 
\end{itemize}
\item {\bf  For} $i=L-1,\cdots,1$
\begin{itemize}

\item ${\bmi R}_{ \uparrow \downarrow }^{i}=
{\bmi S}_{ \uparrow \downarrow }^{i}+{\bmi S}_{ \uparrow \uparrow 
}^{i}
{\bmi  E}_{i+1} {\bmi R}_{ \uparrow \downarrow }^{i+1} {\bmi  E}_{i+1} 
{\left({{\bmi  I}-{\bmi S}_{ \downarrow 
\uparrow }^{i}
{\bmi  E}_{i+1} {\bmi R}_{ \uparrow \downarrow }^{i+1} {\bmi  E}_{i+1} 
}\right)}^{-1}{\bmi S}_{ \downarrow 
\downarrow }^{i}
 $

\item ${\bmi R}_{ \uparrow \uparrow }^{i}= {\bmi S}_{ \uparrow \uparrow 
}^{i}
{\left({{\bmi  I}-{\bmi  E}_{i+1} {\bmi R}_{ \uparrow \downarrow }^{i+1} 
{\bmi 
E}_{i+1} {\bmi S}_{ \downarrow 
\uparrow }^{i}}\right)}^{-1}$

\item ${\bmi R}_{ \downarrow \downarrow }^{i}={\left({{\bmi  I}-{\bmi S}_{ 
\downarrow \uparrow }^{i} {\bmi  E}_{i+1} {\bmi R}_{ \uparrow \downarrow 
}^{i+1} {\bmi  E}_{i+1} }
\right)}^{-1}{\bmi S}_{ \downarrow \downarrow }^{i} $

\item ${\bmi R}_{ \downarrow \uparrow }^{i} = 
{\left({{\bmi  I}-{\bmi S}_{ \downarrow \uparrow }^{i}
{\bmi  E}_{i+1} {\bmi R}_{ \uparrow \downarrow }^{i+1} {\bmi  E}_{i+1} 
}\right)}^{-1}
{\bmi S}_{ \downarrow \uparrow }^{i}$

\item $
\left[{\begin{array}{c}{\bmi \Sigma}^{\rho}_{\uparrow}\\ 
{\bmi \Sigma}^{\rho}_{\downarrow}
\end{array}}\right]=-i\,\left({\begin{array}{cc}
\sqrt{{\bmi E}_{\rho}}&{\bf 0}\\ 
{\bf 0}&\sqrt{{\bmi E}_{\rho}}
\end{array}}\right)\,\hat{\bmi  D}{}^{T}_{z,\rho}(-{\bf 
k}_{x},-{k}_{y})\left({\bmi I}_8\otimes
{\bf DFT}^{T}(-{\bf k}_{x})  \right)\hat{\bmi \Phi}\left[{\CB F}_{\rho}  
\right]
$


\item  $\left[{\begin{array}{c}
{\bmi t}_{ \uparrow }^{i}\\ {\bmi t}_{ \downarrow 
}^{i+1}\end{array}}\right]=
\left[{\begin{array}{cc}
{\bmi R}_{ \uparrow \downarrow }^{i}&{\bmi R}_{ \uparrow \uparrow }^{i}\\ 
{\bmi R}_{ \downarrow \downarrow }^{i}&{\bmi R}_{ \downarrow \uparrow 
}^{i}\end{array}}\right]
\left[{\begin{array}{c}{\bmi \Sigma}^{i}_{\downarrow}\\ 
{\bmi \Sigma}^{i+1}_{\uparrow} + \left({\bmi  E}_{i+1}+{\mit \delta} {\bmi 
S}^{i+1}_{\uparrow\uparrow} \right){  \bmi t}_{ \uparrow }^{i+1}
\end{array}}\right]$
\end{itemize}
\item {\bf  End}

\begin{itemize}
\item  (*{\em top of layer stack\/ }*) 
\end{itemize}
\item ${\bmi R}_{ \downarrow \uparrow }^{0} = 
{\left({{\bmi  I}-{\bmi S}_{ \downarrow \uparrow }^{[f,{0}]}
{\bmi  E}_{1} {\bmi R}_{ \uparrow \downarrow }^{1} {\bmi  E}_{1} 
}\right)}^{-1}
{\bmi S}_{ \downarrow \uparrow }^{[f,{0}]}$


\item  ${\bmi t}_{ \downarrow }^{1}={\bmi R}_{ \downarrow \uparrow }^{0}
({\bmi \Sigma}^{1}_{\uparrow}+{\bmi  E}_{i}{  \bmi t}_{ \uparrow }^{1})$

\item  ${\mit \Delta}{{\bmi V}}_{ \downarrow }^{1}={\bmi t}_{ \downarrow 
}^{1}$

\end{itemize}

\vspace{6 mm}
\end{minipage}
}
\end{figure}
\vfill

\begin{figure}
%\vspace*{10em}
\fbox{
\begin{minipage} {140.0mm} 
{\bf  Algorithm}~{\bf 6.4} {\em continued\/}

\begin{itemize}
\begin{itemize}

\item  (*{\em downsweep\/ }*) 
\end{itemize}
\item {\bf  For} $i=1,\cdots,L-1$
\begin{itemize}
\item  $\left[{\begin{array}{c}
{\mit \Delta}{{\bmi V}}_{ \uparrow }^{i}\\ {\mit \Delta}{{\bmi V}}_{ 
\downarrow }^{i+1}\end{array}}\right]=\left[{\begin{array}{c}{  \bmi R}_{ 
\uparrow \downarrow }^{i}\\ {  \bmi R}_{ \downarrow \downarrow 
}^{i}\end{array}}\right] \left({\bmi  E}_{i}+{\mit \delta} {\bmi 
S}^{i}_{\downarrow\downarrow} \right) {  {\bmi V}}_{ \downarrow 
}^{i}+\left[{\begin{array}{c}{  \bmi t}_{ \uparrow }^{i}\\ {  
\bmi t}_{ \downarrow }^{i+1}\end{array}}\right] $

\item $
{\CB B}_{\rho}=\hat{\bmi \Phi}{}^{-1}\left[\left({\bmi I}_8\otimes 
{\bf DFT}({\bf k}_{x}) \right)\hat{\bmi  D}_{z,\rho}({\bf k}_{x},{k}_{y})
\left({\begin{array}{c}
\sqrt{{\bmi E}_{\rho}}{\mit \Delta}{{\bmi V}}_{ \downarrow }^{\rho}\\ 
\sqrt{{\bmi E}_{\rho}}{\mit \Delta}{{\bmi V}}_{ \uparrow }^{\rho}
\end{array}}\right)\right]
$

\end{itemize}
\item {\bf  End}


\begin{itemize}
\item  (*{\em bottom of layer stack\/ }*) 
\end{itemize}
\item  	$
{\mit \Delta}{{\bmi V}}_{ \uparrow }^{L}=
{\bmi R}_{ \uparrow \downarrow }^{L}
\left({\bmi  E}_{L}+{\mit \delta} {\bmi 
S}^{L}_{\downarrow\downarrow} \right) {{\bmi V}}_{ \downarrow }^{L} +{\bmi 
t}_{ \uparrow }^{L}$

\item $
{\CB B}_{L}=\hat{\bmi \Phi}{}^{-1}\left[\left({\bmi I}_8\otimes 
{\bf DFT}({\bf k}_{x}) \right)\hat{\bmi  D}_{z,L}({\bf k}_{x},{k}_{y})
\left({\begin{array}{c}
\sqrt{{\bmi E}_{L}}{\mit \Delta}{{\bmi V}}_{ \downarrow }^{L}\\ 
\sqrt{{\bmi E}_{L}}{\mit \Delta}{{\bmi V}}_{ \uparrow }^{L}
\end{array}}\right)\right]
$

\end{itemize}

\vspace{6 mm}
\end{minipage}
}
\end{figure}
\vfill

%%%%%%%%%%%%%%%%%%%%%%%%%%%%%%%%%%%%%%%%%%%%%%%%%%%%%%%%%%%%%%%%%%%%%%%%%%%%%%%%%%%%%%%%%%%


In Fig.\ref{preconit} we compare the convergence rates of the 
preconditioned (PGMRESH) and unpreconditioned GMRESH methods when applied 
to the 
$10\%$ model similar to that of \S\ref{sectioncomparison}, but with 
smoother reference properties.  It is  slightly cheaper to 
calculate   $\tilde{\CB P}_F$  using Alg.{\bf 6.4}, 
than $\rf{\CB P}{r,s}$ from (\ref{CBPdef0}) using Alg.{\bf 5.2} and 
Alg.{\bf 5.3},
because no left/right sweeps are required, although 
we must calculate the action ${\mit \delta} {\CB 
S}^{\rho}_{\downarrow\downarrow,\uparrow\uparrow}\cdot {\bmi 
w}, \ 1 \le \rho \le L$ [{\em cf.\/}(\ref{laydSdefD}) \&  
(\ref{laydSdefU})] 
on various field vectors ${\bmi w}$. In summary, we have replaced 
${\CB H}\cdot {\bmi w}$ with the preconditioned operation $\tilde{\bmi 
\Omega}\left({\CB 
H}\cdot {\bmi w}\right)$ that is slightly less than twice as expensive. 

In Fig.\ref{preconit} we observe that the preconditioned problem required 
$18$ iterations while the unpreconditioned problem required $38$ 
iterations to 
secure the same level of convergence. On the basis of this evidence it 
appears that the preconditioned GMRESH mathod is only slightly cheaper that
the  unpreconditioned method. However, the 
principal disadvantage of the GMRES class of  methods relative to the BCG 
based methods of 
\S\ref{sectioncomparison}, is the need to store all iterates and the 
requirement that each new iterate be orthogonalized with respect to all 
the 
preceding ones. These steps require a large amount of storage and 
computation. Consequently,  halving the number of iterations halves the 
amount of re-orthogonalization and storage, which is certainly a 
worthwhile outcome. Also, the BCG methods of \S\ref{sectioncomparison} 
require the application of both  ${\CB H}$ and its dual ${\CB H}^{\dagger}$
at each iteration,  a cost similar to $\tilde{\bmi \Omega}\left({\CB 
H}\cdot {\bmi w}\right)$ in the PGMRES method. In conclusion, on the basis 
of the numerical experiments conducted to date, the {\em 
preconditioned GMRES method is the best overall method for the 
computation single source/multi-receiver problems with strong heterogeneity
superimposed on a reference stratification with large gradients in the 
material properties\/}, at least for the in-plane $P/SV$ problem. 


\begin{figure}
\HideDisplacementBoxes
%\ShowDisplacementBoxes
\centerline{\BoxedEPSF{preconit.EPSF scaled 800}}
\caption{Residual Vs. iteration for $10\%$ heterogeneity using a 
preconditioned generalized minimal residual method}
\protect\label{preconit}
\end{figure}


%%%%%%%%%%%%%%%%%%%%%%%%%%%%%%%%%%%%%%%%%%%%%%%%%%%%%%%%%%%%%%%%%%%%%%%%%%%%%

%%%%%%%%%%%%%%%%%%%%%%%%%%%%%%%%%%%%%%%%%%%%%%%%%%%%%%%%%%%%%%%%%%%%%%%%%%%%%

%%%%%%%%%%%%%%%%%%%%%%%%%%%%%%%%%%%%%%%%%%%%%%%%%%%%%%%%%%%%%%%%%%%%%%%%%%%%%

%%%%%%%%%%%%%%%%%%%%%%%%%%%%%%%%%%%%%%%%%%%%%%%%%%%%%%%%%%%%%%%%%%%%%%%%%%%%%

